\documentclass[12pt, a4paper, oneside]{ctexart}
\usepackage{amsmath, amsthm, amssymb, bm, color, framed, graphicx, mathrsfs}
\usepackage{enumitem} % 用于精细控制列表格式
\usepackage{geometry} % 调整页边距
\usepackage{booktabs} % 三线表
\usepackage{float}    % 浮动体控制
\usepackage{array}    % 数组和表格
\usepackage{mathtools} % 数学工具，提供 dcases 等环境

% 页面设置
\geometry{top=2.5cm, bottom=2.5cm, left=3cm, right=2.5cm}

% 设置列表格式
\setlist[enumerate]{label=\textbf{T\arabic*}, left=0pt, itemindent=*, labelsep=1.5em, align=left, topsep=0.5\baselineskip, partopsep=0pt}
\setlist[enumerate, 2]{label=(\alph*)}

% 自定义环境
\definecolor{shadecolor}{RGB}{241, 241, 255}
\newcounter{problemname}
% 在导言区修改 problem 环境定义
\newenvironment{problem}{
    \stepcounter{problemname}
    \par\noindent
    \textbf{Problem\,\arabic{problemname} }%括号写 题目\arabic{problemname}.
    \par
}{\par}

\newenvironment{solution}{\par\noindent\textbf{ }}{\par}%括号写题目
\newenvironment{note}{\par\noindent\textbf{}}{\par}%括号写注记

% 文档信息
\title{\textbf{物理模型推导题目ver1.3}}
\author{GRS}
\date{2026年4月}
\linespread{1.5}

\begin{document}

% ========== 标题页 ==========
\begin{titlepage}
    \centering
    \vspace*{2cm}
    {\Huge \bfseries 物理模型推导题目} \\[1cm]
    {\Large Theoretical Physics Derivation Problems} \\[2cm]
    {\Large GRS} \\[0.5cm]
    {\Large 2026年3月} \\[2cm]
\end{titlepage}
% ===========================

% ========== 引言 ==========
\vspace{2cm}
\noindent\textbf{}

\vspace{1cm}
\hfill \begin{minipage}{0.8\textwidth}
\centering
\itshape
重要的 200 级公式都要记住。
\end{minipage}
\begin{flushright}
———— Spider
\end{flushright}
\vspace{2cm}
% ===========================

\newpage

% ========== Part 1 力学 ==========
\section*{Part 1 力学}
\begin{enumerate}[label=\textbf{}, leftmargin=*]
    \item \begin{problem}
    证明抛体运动包络线方程:
    \begin{equation*}
    y = \dfrac{v_0^2}{2g} - \dfrac{g}{2v_0^2}x^2.
    \end{equation*}
    \end{problem}
    \begin{solution}
    % 解答留空
    \end{solution}
    \begin{note}
    % 注记留空
    \end{note}

    \item \begin{problem}
    证明曲率半径:
    \begin{equation*}
    \rho = \dfrac{(1+y'^2)^{3/2}}{|y''|} = \dfrac{1}{\kappa}.
    \end{equation*}
    \end{problem}
    \begin{solution}
    \end{solution}
    \begin{note}
    \end{note}

    \item \begin{problem}
    证明软绳有约束:
    \begin{equation*}
    \dfrac{\partial v_c}{\partial s} + \dfrac{v_n}{\rho} = 0.
    \end{equation*}
    \end{problem}
    \begin{solution}
    \end{solution}
    \begin{note}
    \end{note}

    \item \begin{problem}
    证明 $\left(m,M,Wall\right)$ 系统中 $M = k m$ 物块与 $m$ 物块在 $k \gg 1$ 时，碰撞次数:
    \begin{equation*}
    n = \lfloor \sqrt{k} \pi \rfloor.
    \end{equation*}
    \end{problem}
    \begin{solution}
    \end{solution}
    \begin{note}
    \end{note}

    \item \begin{problem}
    导出平方反比力轨道:
    \begin{equation*}
    r = \dfrac{\dfrac{L^2}{GMm^2}}{1+\sqrt{1+\dfrac{2EL^2}{G^2M^2m^3}}\cos\theta} = \dfrac{p}{1+e\cos\theta}.
    \end{equation*}
    \end{problem}
    \begin{solution}
    \end{solution}
    \begin{note}
    \end{note}

    \item \begin{problem}
    导出立方反比力轨道，分 $E>0, e>1$, $E=0, e=1$, $E<0, e<1$ 的排列组合讨论.
    \end{problem}
    \begin{solution}
    \end{solution}
    \begin{note}
    \end{note}

    \item \begin{problem}
    求 $V = -\dfrac{GMm}{r} + \dfrac{1}{2} m \alpha r^2$, $\quad \alpha \ll 1$ 微扰下的进动角速度.
    \end{problem}
    \begin{solution}
    \end{solution}
    \begin{note}
    \end{note}

    \item \begin{problem}
    求 $V = -\dfrac{GMm}{r} + \dfrac{\beta}{r^{1+k}}$ 微扰下的进动角速度.
    \end{problem}
    \begin{solution}
    \end{solution}
    \begin{note}
    \end{note}

    \item \begin{problem}
    长 $l$ 的曲杆，铰接于点 $A(x_0,0,0)$ 处，另一端在 $B(0, y_0, f(y_0))$ 处与钢丝接触。平衡的最小值为:
    \begin{equation*}
    \mu_m = \dfrac{x_0 f'(y_0) - y_0}{x_0 \sqrt{x_0^2 + y_0^2 + f'(y_0)^2}}.
    \end{equation*}
    \end{problem}
    \begin{solution}
    \end{solution}
    \begin{note}
    \end{note}

    \item \begin{problem}
    导出悬链线方程:
    \begin{equation*}
    y = \dfrac{T_0}{\lambda g} \left( \cosh\left( \dfrac{\lambda g}{T_0} x \right) - 1 \right).
    \end{equation*}
    \end{problem}
    \begin{solution}
    \end{solution}
    \begin{note}
    \end{note}
    
    \item \begin{problem}
    导出引力场悬链线:
    \begin{equation*}
    r = -\dfrac{\dfrac{1-A}{A} r_0}{1 - \dfrac{1}{A} \cos(\sqrt{1-A^2} \theta)}.
    \end{equation*}
    \end{problem}
    \begin{solution}
    \end{solution}
    \begin{note}
    \end{note}
    
    \item \begin{problem}
    导出卢瑟福散射在质心系下的微分截面:
    \begin{equation*}
    \dfrac{d\sigma}{d\Omega} = \left( \dfrac{1}{4\pi\epsilon_0} \dfrac{Z_1 Z_2 e^2}{4E} \right)^2 \dfrac{1}{\sin^4(\theta_c/2)}.
    \end{equation*}
    \end{problem}
    \begin{solution}
    \end{solution}
    \begin{note}
    \end{note}
    
    \item \begin{problem}
    导出卢瑟福散射在实验室系下的微分截面:
    \begin{equation*}
    \dfrac{d\sigma}{d\Omega} = \left( \dfrac{1}{4\pi\epsilon_0} \dfrac{Z_1 Z_2 e^2}{2E_1} \right)^2 \dfrac{ \left[ \cos\theta_L + \sqrt{1 - \left( \dfrac{m_1}{m_2} \sin\theta_L \right)^2} \right]^2 }{ \sqrt{1 - \left( \dfrac{m_1}{m_2} \sin\theta_L \right)^2} \sin^4\theta_L }.
    \end{equation*}
    \end{problem}
    \begin{solution}
    \end{solution}
    \begin{note}
    \end{note}
    
    \item \begin{problem}
    导出欧拉运动学方程:
    \begin{equation*}
    \omega_1 = \dot{\phi} \sin\theta \sin\psi + \dot{\theta} \cos\psi,
    \end{equation*}
    \begin{equation*}
    \omega_2 = \dot{\phi} \sin\theta \cos\psi - \dot{\theta} \sin\psi,
    \end{equation*}
    \begin{equation*}
    \omega_3 = \dot{\phi} \cos\theta + \dot{\psi}.
    \end{equation*}
    \end{problem}
    \begin{solution}
    \end{solution}
    \begin{note}
    \end{note}
    
    \item \begin{problem}
    导出欧拉动力学方程:
    \begin{equation*}
    I_1 \dot{\omega}_i - (I_j - I_k) \omega_j \omega_k = M_i, \quad (i,j,k) \text{轮换}.
    \end{equation*}
    \end{problem}
    \begin{solution}
    \end{solution}
    \begin{note}
    \end{note}
    
    \item \begin{problem}
    导出轴对称陀螺在 $\vec{\Pi}$ 空间（角动量空间）的轨道:
    \begin{equation*}
    \Pi_1^2 + \Pi_2^2 + \Pi_3^2 = L^2,
    \end{equation*}
    \begin{equation*}
    \dfrac{\Pi_1^2}{2I_1T} + \dfrac{\Pi_2^2}{2I_2T} + \dfrac{\Pi_3^2}{2I_3T} = 1.
    \end{equation*}
    \end{problem}
    \begin{solution}
    \end{solution}
    \begin{note}
    \end{note}
    
    \item \begin{problem}
    导出自由不对称陀螺的角速度:
    \begin{equation*}
    \Omega_1(\tau) = \sqrt{ \dfrac{2E I_1 - L^2}{I_1(I_3 - I_1)} } \text{cn}(\tau),
    \end{equation*}
    \begin{equation*}
    \Omega_2(\tau) = \sqrt{ \dfrac{2E I_2 - L^2}{I_2(I_3 - I_2)} } \text{sn}(\tau),
    \end{equation*}
    \begin{equation*}
    \Omega_3(\tau) = \sqrt{ \dfrac{L^2 - 2E I_3}{I_3(I_3 - I_1)} } \text{dn}(\tau).
    \end{equation*}
    \end{problem}
    \begin{solution}
    \end{solution}
    \begin{note}
    \end{note}
    
    \item \begin{problem}
    证明对称陀螺定点转动:
    \begin{equation*}
    \dot{\phi} = \dfrac{L_3 - L_2 \cos\theta}{I_1' \sin^2\theta},
    \end{equation*}
    \begin{equation*}
    \dot{\psi} = \dfrac{L_2}{I_3} - \dfrac{L_3 - L_2 \cos\theta}{I_1' \sin^2\theta} \cos\theta,
    \end{equation*}
    \begin{equation*}
    \text{s.t. } I_1' = I_1 + ml^2.
    \end{equation*}
    \end{problem}
    \begin{solution}
    \end{solution}
    \begin{note}
    \end{note}
    
    \item \begin{problem}
    证明刚体所有运动构成的流形 $Q$ 有:
    \begin{equation*}
    \mathbb{Q} = \mathbb{R}^3 \rtimes SO(3).
    \end{equation*}
    \end{problem}
    \begin{solution}
    \end{solution}
    \begin{note}
    \end{note}
    
    \item \begin{problem}
    找到天体运动拉格朗日点 $L_i, i=1,2,3,4,5$，并讨论运动模式.
    \end{problem}
    \begin{solution}
    \end{solution}
    \begin{note}
    \end{note}
    
    \item \begin{problem}
    导出平面双振子在 $\phi_1(0) = -\phi_2(0) = \phi_0 \ll 1$, $\dot{\phi}_1(0)=\dot{\phi}_2(0)=0$ 初态下系统的解
    \begin{equation*}
    \phi_1(t) = \dfrac{\phi_0}{4} \begin{pmatrix}
        (2+\sqrt{2})\cos(\omega_1 t) + (2-\sqrt{2})\cos(\omega_2 t) \\
        -(2+2\sqrt{2})\cos(\omega_1 t) + (2\sqrt{2}-2)\cos(\omega_2 t)
    \end{pmatrix},
    \end{equation*}
    \begin{equation*}
    \text{s.t. } \omega_{1,2} = \sqrt{\dfrac{k}{m}(2 \pm \sqrt{2})}.
    \end{equation*}
    \end{problem}
    \begin{solution}
    \end{solution}
    \begin{note}
    \end{note}
    
    \item \begin{problem}
    导出成环原子链的本征频率:
    \begin{equation*}
    \omega^2(p) = 2\omega_0^2 (1-\cos(p)),
    \end{equation*}
    \begin{equation*}
    p = \dfrac{2\pi l}{N}, \quad l \in \left[-\dfrac{N}{2}, \dfrac{N}{2}\right].
    \end{equation*}
    \end{problem}
    \begin{solution}
    \end{solution}
    \begin{note}
    \end{note}
    
    \item \begin{problem}
    导出单摆周期的修正:
    \begin{equation*}
    T = 2\pi \sqrt{\dfrac{l}{g}} \left( 1 + \dfrac{1}{16} \theta_0^2 + \cdots \right).
    \end{equation*}
    \end{problem}
    \begin{solution}
    \end{solution}
    \begin{note}
    \end{note}
    
    \item \begin{problem}
    对参数振子: $\omega^2(t) = \omega_0^2 \bigl( 1 + h \cos(\gamma t) \bigr)$, $h \ll 1$, $\gamma = 2\omega_0 + 2\varepsilon$, $\varepsilon \ll \omega_0$，存在指数增长解的条件:
    \begin{equation*}
    |\varepsilon| < \dfrac{1}{4} h \omega_0.
    \end{equation*}
    \end{problem}
    \begin{solution}
    \end{solution}
    \begin{note}
    \end{note}
    
    \item \begin{problem}
    导出二维欧氏平面上在 $V = \dfrac{1}{2} k (x^2 + y^2)$ 与 $\vec{F} = -m \vec{\Omega} \times \vec{v}$ 下质点 $m$ 的运动模式.
    \end{problem}
    \begin{solution}
    \end{solution}
    \begin{note}
    \end{note}
    
    \item \begin{problem}
    导出一维双原子链 $m, M$ 的色散关系:
    \begin{equation*}
    \omega_{\pm}^2 = \dfrac{k(M+m)}{Mm} \left\{ 1 \pm \sqrt{ 1 - \dfrac{4Mm \sin^2 (\frac{1}{2}ka)}{(M+m)^2} } \right\}.
    \end{equation*}
    \end{problem}
    \begin{solution}
    \end{solution}
    \begin{note}
    \end{note}
    
    \item \begin{problem}
    导出 N-C 方程:
    \begin{equation*}
    \rho \dfrac{d^2 \vec{\xi}}{dt^2} = (\lambda + \mu) \nabla (\nabla \cdot \vec{\xi}) + \mu \nabla^2 \vec{\xi}.
    \end{equation*}
    \end{problem}
    \begin{solution}
    \end{solution}
    \begin{note}
    \end{note}
    
    \item \begin{problem}
    导出弹性波波速:
    \begin{equation*}
    c_{//} = \sqrt{\dfrac{\lambda + 2\mu}{\rho}}, \quad c_{\bot} = \sqrt{\dfrac{\mu}{\rho}}.
    \end{equation*}
    \end{problem}
    \begin{solution}
    \end{solution}
    \begin{note}
    \end{note}
    
    \item \begin{problem}
    证明:
    \begin{equation*}
    \nu = \dfrac{\lambda}{2(\lambda+\mu)},
    \end{equation*}
    \begin{equation*}
    E = \dfrac{\mu(3\lambda+2\mu)}{\lambda+\mu},
    \end{equation*}
    \begin{equation*}
    \lambda = \dfrac{E\nu}{(1+\nu)(1-2\nu)},
    \end{equation*}
    \begin{equation*}
    \mu = \dfrac{E}{2(1+\nu)}.
    \end{equation*}
    \end{problem}
    \begin{solution}
    \end{solution}
    \begin{note}
    \end{note}
    
    \item \begin{problem}
    导出两边固定、中间拱起的轻绳形状近似解:
    \begin{equation*}
    y = \dfrac{L \Delta}{\pi} \left( 1 - \cos\left(\dfrac{2\pi x}{L}\right) \right).
    \end{equation*}
    \end{problem}
    \begin{solution}
    \end{solution}
    \begin{note}
    \end{note}
    
    \item \begin{problem}
    质量 $m$ 的弹簧振子以 $\omega$ 频率振动时的机械能:
    \begin{equation*}
    x = \dfrac{ \sin\left( \sqrt{\dfrac{m}{k}} \dfrac{\omega}{l} \xi \right) }{ \sqrt{\dfrac{m}{k}} \dfrac{\omega}{l} \cos\left( \sqrt{\dfrac{m}{k}} \dfrac{\omega}{l} \right) }.
    \end{equation*}
    \end{problem}
    \begin{solution}
    \end{solution}
    \begin{note}
    \end{note}
    
    \item \begin{problem}
    无原长 $l \ll \dfrac{mg}{k}$ 质量 $m$ 弹簧悬链线:
    \begin{equation*}
    y = -\dfrac{mg}{k} \left( \dfrac{x}{d} \right)^2 + \dfrac{mg}{k} \left( \dfrac{x}{d} \right).
    \end{equation*}
    \end{problem}
    \begin{solution}
    \end{solution}
    \begin{note}
    \end{note}
    
    \item \begin{problem}
    导出泊肃叶公式:
    \begin{equation*}
    Q = \dfrac{\pi R^4 \Delta p}{8\eta L}.
    \end{equation*}
    \end{problem}
    \begin{solution}
    \end{solution}
    \begin{note}
    \end{note}
    
    \item \begin{problem}
    证明欧拉动力学方程（流体）:
    \begin{equation*}
    \dfrac{\partial}{\partial t} (\rho \vec{v}) + \nabla \cdot (\rho \vec{v} \otimes \vec{v}) + \nabla p = 0.
    \end{equation*}
    \end{problem}
    \begin{solution}
    \end{solution}
    \begin{note}
    \end{note}
    
    \item \begin{problem}
    证明理想流体，准静态绝热下能量守恒方程:
    \begin{equation*}
    \dfrac{\partial}{\partial t} \left( u_s + \dfrac{v^2}{2} \right) + \nabla \cdot \left[ \left( u_s + \dfrac{v^2}{2} + \dfrac{p}{\rho} \right) \vec{v} \right] = 0.
    \end{equation*}
    \end{problem}
    \begin{solution}
    \end{solution}
    \begin{note}
    \end{note}
    
    \item \begin{problem}
    证明 N-S 方程:
    \begin{equation*}
    \rho \left( \dfrac{\partial \vec{v}}{\partial t} + (\vec{v} \cdot \nabla) \vec{v} \right) = -\nabla p + \eta \nabla^2 \vec{v} + \vec{f}.
    \end{equation*}
    \end{problem}
    \begin{solution}
    \end{solution}
    \begin{note}
    \end{note}
    
    \item \begin{problem}
    讨论阻尼受迫振动的运动模式.
    \end{problem}
    \begin{solution}
    \end{solution}
    \begin{note}
    \end{note}
    
    \item \begin{problem}
    证明经典多普勒公式:
    \begin{equation*}
    \nu = \nu_0 \dfrac{u - v_B}{u - v_S \cos\varphi_S}.
    \end{equation*}
    \end{problem}
    \begin{solution}
    \end{solution}
    \begin{note}
    \end{note}
    
    \item \begin{problem}
    单摆，悬挂点受 $\vec{s} = \delta_0 (\cos\alpha \, \vec{\imath} + \sin\alpha \, \vec{\jmath}) \cos(\omega t)$ 驱动，证明微扰下低频小振动频率:
    \begin{equation*}
    \Omega \approx \sqrt{ \dfrac{g}{l} \cos\theta - \dfrac{\delta_0^2 \cos(2\theta-2\alpha) \omega^2}{2l^2} }.
    \end{equation*}
    \end{problem}
    \begin{solution}
    \end{solution}
    \begin{note}
    \end{note}
    
    \item \begin{problem}
    $m_1, m_2, l_1, l_2$ 双摆，驱动 $\delta = \delta_0 \cos(\omega t)$ 求简正振动频率满足的方程: $A\omega^4 + B\omega^2 + C = 0$.
    其中:
    \begin{equation*}
    A = 1,
    \end{equation*}
    \begin{equation*}
    B = -\dfrac{m_1+m_2}{m_1} \left[ g \left( \dfrac{1}{l_1} + \dfrac{1}{l_2} \right) + \dfrac{(m_1+m_2)^2 (l_1+l_2)}{2 m_1 l_1^2 l_2^2} \delta_0^2 \omega^2 + \dfrac{m_2 \delta_0^2 \omega^2}{m_1 l_1 l_2} \right],
    \end{equation*}
    \begin{equation*}
    C = \dfrac{m_1+m_2}{m_1 l_1 l_2} g^2 + \dfrac{(m_1+m_2)^2}{2m_1^2 l_1^2 l_2^2} (l_1+l_2) \delta_0^2 \omega^2 g + \dfrac{(m_1+m_2)^2}{4 m_1^2 l_1^2 l_2^2} \delta_0^4 \omega^4.
    \end{equation*}
    \end{problem}
    \begin{solution}
    \end{solution}
    \begin{note}
    \end{note}
\end{enumerate}
% ========== Part 1 结束 ==========

\newpage
% ========== Part 2 热学 ==========
\section*{Part 2 热学}
\begin{enumerate}[label=\textbf{}, leftmargin=*]
    \item \begin{problem}
    证明 Maxwell 关系:
    \begin{equation*}
    \left( \dfrac{\partial T}{\partial V} \right)_S = -\left( \dfrac{\partial p}{\partial S} \right)_V,
    \end{equation*}
    \begin{equation*}
    \left( \dfrac{\partial T}{\partial p} \right)_S = \left( \dfrac{\partial V}{\partial S} \right)_p,
    \end{equation*}
    \begin{equation*}
    \left( \dfrac{\partial S}{\partial V} \right)_T = \left( \dfrac{\partial p}{\partial T} \right)_V,
    \end{equation*}
    \begin{equation*}
    \left( \dfrac{\partial S}{\partial p} \right)_T = -\left( \dfrac{\partial V}{\partial T} \right)_p.
    \end{equation*}
    \end{problem}
    \begin{solution}
    \end{solution}
    \begin{note}
    \end{note}

    \item \begin{problem}
    证明能态、焓态定理:
    \begin{equation*}
    \left( \dfrac{\partial U}{\partial V} \right)_T = T \left( \dfrac{\partial p}{\partial T} \right)_V - p,
    \end{equation*}
    \begin{equation*}
    \left( \dfrac{\partial H}{\partial p} \right)_T = T \left( \dfrac{\partial V}{\partial T} \right)_p + V.
    \end{equation*}
    \end{problem}
    \begin{solution}
    \end{solution}
    \begin{note}
    \end{note}

    \item \begin{problem}
    证明响应函数间有:
    \begin{equation*}
    \alpha = \beta \kappa_T p,
    \end{equation*}
    \begin{equation*}
    \dfrac{C_p}{C_V} = \dfrac{\kappa_T}{\kappa_S},
    \end{equation*}
    \begin{equation*}
    C_p - C_V = TV\alpha\beta.
    \end{equation*}
    \end{problem}
    \begin{solution}
    \end{solution}
    \begin{note}
    \end{note}

    \item \begin{problem}
    证明 Gibbs-Helmholtz 方程:
    \begin{equation*}
    \left( \dfrac{\partial (G/T)}{\partial T} \right)_p = -\dfrac{H}{T^2}.
    \end{equation*}
    \end{problem}
    \begin{solution}
    \end{solution}
    \begin{note}
    \end{note}

    \item \begin{problem}
    证明绝热大气满足:
    \begin{equation*}
    n = n_0 \left( 1 - (1-\dfrac{1}{\gamma}) \dfrac{mgz}{kT_0} \right)^{\frac{1}{\gamma-1}}.
    \end{equation*}
    \end{problem}
    \begin{solution}
    \end{solution}
    \begin{note}
    \end{note}

    \item \begin{problem}
    证明等温大气满足:
    \begin{equation*}
    n = n_0 \exp \left( -\dfrac{m g h}{k T} \right).
    \end{equation*}
    \end{problem}
    \begin{solution}
    \end{solution}
    \begin{note}
    \end{note}

    \item \begin{problem}
    证明范德瓦尔斯气体节流的 J-T 系数:
    \begin{equation*}
    \mu = \left( \dfrac{\partial T}{\partial p} \right)_H = \dfrac{2aV_m(V_m-b)^2 - bRT V_m^3}{R^2 T V_m^3 + C_{V,m} [RTV_m^3 - 2a(V_m-b)^2]}.
    \end{equation*}
    \end{problem}
    \begin{solution}
    \end{solution}
    \begin{note}
    \end{note}

    \item \begin{problem}
    证明声速公式:
    \begin{equation*}
    c = \sqrt{\gamma \dfrac{RT}{M}}.
    \end{equation*}
    \end{problem}
    \begin{solution}
    \end{solution}
    \begin{note}
    \end{note}

    \item \begin{problem}
    证明奥托循环效率:
    \begin{equation*}
    \eta_{\text{Otto}} = 1 - \dfrac{1}{r^{\gamma-1}}, \quad r = \dfrac{V_1}{V_2}.
    \end{equation*}
    \end{problem}
    \begin{solution}
    \end{solution}
    \begin{note}
    \end{note}

    \item \begin{problem}
    证明狄塞尔循环效率:
    \begin{equation*}
    \eta_{\text{Diesel}} = 1 - \dfrac{\rho^{\gamma} - 1}{\gamma(\rho-1) r^{\gamma-1}}, \quad r = \dfrac{V_1}{V_2}, \quad \rho = \dfrac{V_3}{V_2}.
    \end{equation*}
    \end{problem}
    \begin{solution}
    \end{solution}
    \begin{note}
    \end{note}

    \item \begin{problem}
    证明范德瓦尔斯气体的内能和焓:
    \begin{equation*}
    U = \int_{T_0}^{T} C_V dT - \dfrac{\nu^2 a}{V} + U_0,
    \end{equation*}
    \begin{equation*}
    H = \int_{T_0}^{T} C_V dT + \left( \dfrac{\nu RTV}{V-\nu b} - \dfrac{2\nu^2 a}{V} \right) + H_0.
    \end{equation*}
    \end{problem}
    \begin{solution}
    \end{solution}
    \begin{note}
    \end{note}

    \item \begin{problem}
    证明张力系统中:
    \begin{equation*}
    S_S = -A_S \left( \dfrac{\partial \sigma}{\partial T} \right)_{A_S}, \quad U_S = \sigma A_S - T A_S \left( \dfrac{\partial \sigma}{\partial T} \right)_{A_S}.
    \end{equation*}
    \end{problem}
    \begin{solution}
    \end{solution}
    \begin{note}
    \end{note}

    \item \begin{problem}
    证明附加压强的 Laplace 公式:
    \begin{equation*}
    \Delta p = \sigma \left( \dfrac{1}{R_1} + \dfrac{1}{R_2} \right).
    \end{equation*}
    \end{problem}
    \begin{solution}
    \end{solution}
    \begin{note}
    \end{note}

    \item \begin{problem}
    证明方形容器附加高度:
    \begin{equation*}
    h = \sqrt{ \dfrac{2\sigma}{\rho g} (1 - \sin\theta) }.
    \end{equation*}
    \end{problem}
    \begin{solution}
    \end{solution}
    \begin{note}
    \end{note}

    \item \begin{problem}
    证明范德瓦尔斯气体临界参数:
    \begin{equation*}
    T_c = \dfrac{8a}{27Rb},
    \end{equation*}
    \begin{equation*}
    V_c = 3b,
    \end{equation*}
    \begin{equation*}
    p_c = \dfrac{a}{27b^2}.
    \end{equation*}
    \end{problem}
    \begin{solution}
    \end{solution}
    \begin{note}
    \end{note}

    \item \begin{problem}
    导出饱和蒸气压方程:
    \begin{equation*}
    \ln \left( \dfrac{P_S}{P_{S,0}} \right) = B \left( 1 - \dfrac{T_0}{T} \right) + C \ln \left( \dfrac{T}{T_0} \right).
    \end{equation*}
    \end{problem}
    \begin{solution}
    \end{solution}
    \begin{note}
    \end{note}

    \item \begin{problem}
    导出诱导核的中肯半径:
    \begin{equation*}
    r_C = \dfrac{2\sigma \rho_G}{(\rho_L - \rho_G)(P' - P)}.
    \end{equation*}
    \end{problem}
    \begin{solution}
    \end{solution}
    \begin{note}
    \end{note}

    \item \begin{problem}
    导出 Van't Hoff 方程:
    \begin{equation*}
    \dfrac{d \ln K_P(T)}{dT} = \dfrac{\Delta H}{RT^2}.
    \end{equation*}
    \end{problem}
    \begin{solution}
    \end{solution}
    \begin{note}
    \end{note}

    \item \begin{problem}
    导出正则系综的能量涨落:
    \begin{equation*}
    \langle (\Delta E)^2 \rangle = kT^2 C_V.
    \end{equation*}
    \end{problem}
    \begin{solution}
    \end{solution}
    \begin{note}
    \end{note}

    \item \begin{problem}
    导出巨正则系综的粒子数涨落和能量涨落:
    \begin{equation*}
    \langle (\Delta N)^2 \rangle = \dfrac{kT \bar{N}^2}{V} \kappa_T,
    \end{equation*}
    \begin{equation*}
    \langle (\Delta E)^2 \rangle = kT^2 C_V + \left( \dfrac{\partial \bar{E}}{\partial \bar{N}} \right)_{T,V}^2 \dfrac{kT \bar{N}^2}{V} \kappa_T.
    \end{equation*}
    \end{problem}
    \begin{solution}
    \end{solution}
    \begin{note}
    \end{note}

    \item \begin{problem}
    导出单原子分子理想气体的熵:
    \begin{equation*}
    S = Nk \ln \left[ \dfrac{V}{N} \left( \dfrac{2\pi m kT}{h^2} \right)^{3/2} \right] + \dfrac{5}{2} Nk.
    \end{equation*}
    \end{problem}
    \begin{solution}
    \end{solution}
    \begin{note}
    \end{note}

    \item \begin{problem}
    导出双原子分子气体 $H_2$ 的定容热容 $C_V$.
    \end{problem}
    \begin{solution}
    \end{solution}
    \begin{note}
    \end{note}

    \item \begin{problem}
    导出电场中双原子分子气体的电极化率:
    \begin{equation*}
    P_0 = \dfrac{N d_0}{V} \left( \coth(\beta d_0 E_0) - \dfrac{1}{\beta d_0 E_0} \right) E.
    \end{equation*}
    \end{problem}
    \begin{solution}
    \end{solution}
    \begin{note}
    \end{note}

    \item \begin{problem}
    导出混合理想气体的熵和化学势:
    \begin{equation*}
    S = \sum_i N_i k \ln \left[ \dfrac{V}{N_i \lambda_i^3} I_i(T) \right] + N_i k T \dfrac{d \ln I_i(T)}{dT} + \dfrac{5}{2} N_i k,
    \end{equation*}
    \begin{equation*}
    \mu_i = \mu_i^\Theta(T) + kT \ln \dfrac{p_i}{p^\Theta}.
    \end{equation*}
    \end{problem}
    \begin{solution}
    \end{solution}
    \begin{note}
    \end{note}

    \item \begin{problem}
    导出光子气:
    \begin{equation*}
    u(\omega, T) = \dfrac{\hbar}{\pi^2 c^3} \dfrac{\omega^3}{e^{\hbar\omega/kT}-1},
    \end{equation*}
    \begin{equation*}
    E = \sigma V T^4,
    \end{equation*}
    \begin{equation*}
    P = \dfrac{\sigma}{3} T^4,
    \end{equation*}
    \begin{equation*}
    G = 0,
    \end{equation*}
    \begin{equation*}
    S = \dfrac{4}{3} \sigma V T^3.
    \end{equation*}
    \end{problem}
    \begin{solution}
    \end{solution}
    \begin{note}
    \end{note}

    \item \begin{problem}
    证明爱因斯坦热容:
    \begin{equation*}
    C_V = 3Nk \dfrac{(\beta \hbar \omega)^2 e^{\beta \hbar \omega}}{(e^{\beta \hbar \omega} - 1)^2}.
    \end{equation*}
    \end{problem}
    \begin{solution}
    \end{solution}
    \begin{note}
    \end{note}

    \item \begin{problem}
    证明德拜低温热容:
    \begin{equation*}
    C_V = \dfrac{12\pi^4}{5} Nk \left( \dfrac{T}{\theta_D} \right)^3, \quad \text{s.t. } \theta_D = \dfrac{\hbar \omega_D}{k}.
    \end{equation*}
    \end{problem}
    \begin{solution}
    \end{solution}
    \begin{note}
    \end{note}

    \item \begin{problem}
    导出居里定律:
    \begin{equation*}
    \chi = \dfrac{C}{T}.
    \end{equation*}
    \end{problem}
    \begin{solution}
    \end{solution}
    \begin{note}
    \end{note}

    \item \begin{problem}
    导出二能级系统热容:
    \begin{equation*}
    C_V = Nk \left( \dfrac{\Delta}{kT} \right)^2 \dfrac{e^{\Delta/kT}}{(1+e^{\Delta/kT})^2}, \quad \Delta = 2\epsilon.
    \end{equation*}
    \end{problem}
    \begin{solution}
    \end{solution}
    \begin{note}
    \end{note}

    \item \begin{problem}
    证明弱简并理想玻色气体状态方程:
    \begin{equation*}
    \dfrac{PV}{NkT} = 1 - \dfrac{\zeta(5/2)}{2^{5/2}} \dfrac{\lambda^3}{v} - \cdots, \quad \text{s.t. } \lambda = \sqrt{\dfrac{h^2}{2\pi m kT}}.
    \end{equation*}
    \end{problem}
    \begin{solution}
    \end{solution}
    \begin{note}
    \end{note}

    \item \begin{problem}
    证明弱简并理想玻色气体的熵:
    \begin{equation*}
    S = \dfrac{5}{2} Nk \dfrac{g_{5/2}(z)}{g_{3/2}(z)} - Nk \ln z.
    \end{equation*}
    \end{problem}
    \begin{solution}
    \end{solution}
    \begin{note}
    \end{note}

    \item \begin{problem}
    验证 BEC 热容偏导在 $T_c$ 处不连续.
    \end{problem}
    \begin{solution}
    \end{solution}
    \begin{note}
    \end{note}

    \item \begin{problem}
    证明 $V = \dfrac{1}{2} m (\omega_1^2 x^2 + \omega_2^2 y^2 + \omega_3^2 z^2)$ 中的 BEC 凝聚温度:
    \begin{equation*}
    T_c = \dfrac{\hbar \bar{\omega} N^{1/3}}{k [\zeta(3)]^{1/3}}.
    \end{equation*}
    \end{problem}
    \begin{solution}
    \end{solution}
    \begin{note}
    \end{note}

    \item \begin{problem}
    证明弱简并理想费米气体状态方程:
    \begin{equation*}
    \dfrac{PV}{NkT} = 1 + \dfrac{\zeta(5/2)}{2^{5/2}} \dfrac{\lambda^3}{v} + \cdots.
    \end{equation*}
    \end{problem}
    \begin{solution}
    \end{solution}
    \begin{note}
    \end{note}

    \item \begin{problem}
    证明零温理想费米气体:
    \begin{equation*}
    \epsilon_F = \dfrac{\hbar^2}{2m} (3\pi^2 n)^{2/3},
    \end{equation*}
    \begin{equation*}
    p_F = \hbar (3\pi^2 n)^{1/3},
    \end{equation*}
    \begin{equation*}
    E_0 = \dfrac{3}{5} N \epsilon_F,
    \end{equation*}
    \begin{equation*}
    P_0 = \dfrac{2}{5} n \epsilon_F.
    \end{equation*}
    \end{problem}
    \begin{solution}
    \end{solution}
    \begin{note}
    \end{note}

    \item \begin{problem}
    证明低温理想费米气体热容:
    \begin{equation*}
    C_V = \dfrac{\pi^2}{2} Nk \dfrac{T}{T_F}, \quad \text{s.t. } T_F = \dfrac{\epsilon_F}{k}.
    \end{equation*}
    \end{problem}
    \begin{solution}
    \end{solution}
    \begin{note}
    \end{note}

    \item \begin{problem}
    证明磁场中强简并费米气体的顺磁化率:
    \begin{equation*}
    \chi_P = \dfrac{3}{2} \dfrac{n \mu_B^2}{\epsilon_F}.
    \end{equation*}
    \end{problem}
    \begin{solution}
    \end{solution}
    \begin{note}
    \end{note}

    \item \begin{problem}
    证明 $V = \dfrac{1}{2} m (\omega_1^2 x^2 + \omega_2^2 y^2 + \omega_3^2 z^2)$ 中的强简并费米气体热容:
    \begin{equation*}
    C_V = \pi^2 Nk \left( \dfrac{kT}{\mu_0} \right) + \cdots.
    \end{equation*}
    \end{problem}
    \begin{solution}
    \end{solution}
    \begin{note}
    \end{note}

    \item \begin{problem}
    导出平均泄流数率:
    \begin{equation*}
    \Gamma = \dfrac{1}{4} n \bar{v}.
    \end{equation*}
    \end{problem}
    \begin{solution}
    \end{solution}
    \begin{note}
    \end{note}

    \item \begin{problem}
    导出 $t$ 维 MB 分布:
    \begin{equation*}
    f(\vec{v}) = \left( \dfrac{m}{2\pi kT} \right)^{t/2} e^{-\frac{m v^2}{2kT}},
    \end{equation*}
    \begin{equation*}
    f(v, t) = \dfrac{t \pi^{t/2}}{\Gamma(\frac{t}{2}+1)} v^{t-1} \left( \dfrac{m}{2\pi kT} \right)^{t/2} e^{-\frac{m v^2}{2kT}}.
    \end{equation*}
    \end{problem}
    \begin{solution}
    \end{solution}
    \begin{note}
    \end{note}

    \item \begin{problem}
    证明泄流升维:
    \begin{equation*}
    f(v_e, t) = f(v, t+1).
    \end{equation*}
    \end{problem}
    \begin{solution}
    \end{solution}
    \begin{note}
    \end{note}

    \item \begin{problem}
    导出自由程分布:
    \begin{equation*}
    P(\lambda) = e^{-\lambda / \bar{\lambda}}.
    \end{equation*}
    \end{problem}
    \begin{solution}
    \end{solution}
    \begin{note}
    \end{note}

    \item \begin{problem}
    导出稀薄混合气体平均自由程:
    \begin{equation*}
    \lambda_1 = \dfrac{1}{\sqrt{2} \pi d_1^2 n_1 + \pi (d_1+d_2)^2 n_2 \sqrt{\dfrac{m_1}{16\mu}}}, \quad \lambda_2 = \dfrac{1}{\sqrt{2} \pi d_2^2 n_2 + \pi (d_1+d_2)^2 n_1 \sqrt{\dfrac{m_2}{16\mu}}}.
    \end{equation*}
    \end{problem}
    \begin{solution}
    \end{solution}
    \begin{note}
    \end{note}

    \item \begin{problem}
    维里展开得到实际气体状态方程:
    \begin{equation*}
    \dfrac{p}{kT} = n + B_2(T) n^2 + B_3(T) n^3 + \cdots.
    \end{equation*}
    \end{problem}
    \begin{solution}
    \end{solution}
    \begin{note}
    \end{note}

    \item \begin{problem}
    在 L-J 势下得到范德瓦尔斯方程:
    \begin{equation*}
    p = \dfrac{NkT}{V-b} - \dfrac{a}{V^2}.
    \end{equation*}
    \end{problem}
    \begin{solution}
    \end{solution}
    \begin{note}
    \end{note}

    \item \begin{problem}
    证明物态方程:
    \begin{equation*}
    pV = NkT \left[ 1 - \dfrac{n}{6kT} \int d^3\vec{r} \, (\vec{r} \cdot \nabla U(\vec{r})) g(\vec{r}) \right].
    \end{equation*}
    \end{problem}
    \begin{solution}
    \end{solution}
    \begin{note}
    \end{note}

    \item \begin{problem}
    导出稀薄等离子体的热态方程:
    \begin{equation*}
    p = \dfrac{NkT}{V} - \dfrac{e^2}{3V^{3/2}} \left( \dfrac{\pi}{kT} \right)^{1/2} \left( \sum_i N_i z_i^2 \right)^{3/2}.
    \end{equation*}
    \end{problem}
    \begin{solution}
    \end{solution}
    \begin{note}
    \end{note}

    \item \begin{problem}
    导出平均场近似 Ising 模型临界点热容:
    \begin{equation*}
    C_H = \begin{cases} 
    3Nk/2, & T \to T_c^- \\
    0, & T \to T_c^+ 
    \end{cases}.
    \end{equation*}
    \end{problem}
    \begin{solution}
    \end{solution}
    \begin{note}
    \end{note}

    \item \begin{problem}
    导出一维 Ising 模型配分函数:
    \begin{equation*}
    Z = \left[ e^{\beta J} \cosh(\beta H) + \sqrt{e^{2\beta J} \sinh^2(\beta H) + e^{-2\beta J}} \right]^N, \quad m = \dfrac{\sinh(\beta H)}{\sqrt{\sinh^2(\beta H) + e^{-4\beta J}}}.
    \end{equation*}
    \end{problem}
    \begin{solution}
    \end{solution}
    \begin{note}
    \end{note}

    \item \begin{problem}
    导出海森堡模型临界温度:
    \begin{equation*}
    T_c = \dfrac{qJ}{3k}.
    \end{equation*}
    \end{problem}
    \begin{solution}
    \end{solution}
    \begin{note}
    \end{note}

    \item \begin{problem}
    导出玻尔兹曼积分-微分方程:
    \begin{equation*}
    \dfrac{\partial f}{\partial t} + \nabla_r \cdot (f \dot{\vec{r}}) + \nabla_v \cdot (f \dot{\vec{v}}) = \int (f_1 f' - f_1 f) d\vec{v}_1 \Lambda d\Omega.
    \end{equation*}
    \end{problem}
    \begin{solution}
    \end{solution}
    \begin{note}
    \end{note}

    \item \begin{problem}
    导出输运系数:
    \begin{equation*}
    \eta = nkT \bar{\tau}, \quad \sigma = \dfrac{ne^2 \tau_F}{m}, \quad \kappa = \dfrac{\pi}{3} \dfrac{n \tau_F}{m} kT^2.
    \end{equation*}
    \end{problem}
    \begin{solution}
    \end{solution}
    \begin{note}
    \end{note}

    \item \begin{problem}
    导出非平衡热机最优化功率时的效率:
    \begin{equation*}
    \eta = 1 - \sqrt{\dfrac{T_L}{T_h}}.
    \end{equation*}
    \end{problem}
    \begin{solution}
    \end{solution}
    \begin{note}
    \end{note}

    \item \begin{problem}
    证明爱因斯坦扩散系数:
    \begin{equation*}
    D = \dfrac{kT}{6\pi a \eta}.
    \end{equation*}
    \end{problem}
    \begin{solution}
    \end{solution}
    \begin{note}
    \end{note}
\end{enumerate}
% ========== Part 2 结束 ==========
\newpage
% ========== Part 3 电磁学 ==========
\section*{Part 3 电磁学}
\begin{enumerate}[label=\textbf{}, leftmargin=*, resume]
    \item \begin{problem}
    导出场点处的电偶极矩电场:
    \begin{equation*}
    \vec{E} = \dfrac{1}{4\pi\epsilon_0} \dfrac{3(\hat{R} \cdot \vec{p})\hat{R} - \vec{p}}{R^3}.
    \end{equation*}
    \end{problem}
    \begin{solution}
    \end{solution}
    \begin{note}
    \end{note}

    \item \begin{problem}
    导出电偶极层的电场:
    \begin{equation*}
    \vec{E} = \dfrac{\tau \nabla \Omega}{4\pi \varepsilon_0}.
    \end{equation*}
    \end{problem}
    \begin{solution}
    \end{solution}
    \begin{note}
    \end{note}

    \item \begin{problem}
    导出二维电偶极矩的电场:
    \begin{equation*}
    \vec{E} = \dfrac{1}{2\pi\epsilon_0} \dfrac{2(\hat{R} \cdot \vec{p})\hat{R} - \vec{p}}{R^3}.
    \end{equation*}
    \end{problem}
    \begin{solution}
    \end{solution}
    \begin{note}
    \end{note}

    \item \begin{problem}
    二维电偶极层的电场:
    \begin{equation*}
    \vec{E} = \dfrac{\tau \nabla \theta}{2\pi \varepsilon_0}.
    \end{equation*}
    \end{problem}
    \begin{solution}
    \end{solution}
    \begin{note}
    \end{note}

    \item \begin{problem}
    导出有限长带电细棒的电场:
    \begin{equation*}
    \vec{E} = \dfrac{\lambda}{2\pi\epsilon_0 \sinh \mu} \dfrac{\hat{\mu}}{c \sqrt{\sinh^2\mu + \sin^2\nu}},
    \end{equation*}
    s.t. $x = c \cosh\mu \cos\nu, \quad y = c \sinh\mu \sin\nu$.
    \end{problem}
    \begin{solution}
    \end{solution}
    \begin{note}
    \end{note}

    \item \begin{problem}
    证明点电荷对导体球的电像:
    \begin{equation*}
    q' = -\dfrac{R}{r} q, \quad r' = \dfrac{R^2}{r}.
    \end{equation*}
    \end{problem}
    \begin{solution}
    \end{solution}
    \begin{note}
    \end{note}

    \item \begin{problem}
    证明线电荷对导体圆柱的电像:
    \begin{equation*}
    \lambda' = -\lambda, \quad r' = \dfrac{R^2}{d}.
    \end{equation*}
    \end{problem}
    \begin{solution}
    \end{solution}
    \begin{note}
    \end{note}

    \item \begin{problem}
    证明不带电导体球在均匀外场中的偶极矩:
    \begin{equation*}
    \vec{p} = 4\pi\epsilon_0 R^3 \vec{E}_0.
    \end{equation*}
    \end{problem}
    \begin{solution}
    \end{solution}
    \begin{note}
    \end{note}

    \item \begin{problem}
    导出旋转椭球体电容:
    \begin{equation*}
    C = \begin{cases} 
    \dfrac{4\pi\epsilon_0 \sqrt{a^2-b^2}}{\operatorname{arccosh}(a/b)}, & a>b \\[10pt]
    \dfrac{4\pi\epsilon_0 \sqrt{b^2-a^2}}{\arccos(a/b)}, & a<b 
    \end{cases}.
    \end{equation*}
    \end{problem}
    \begin{solution}
    \end{solution}
    \begin{note}
    \end{note}

    \item \begin{problem}
    导出两相同导体球电容:
    \begin{equation*}
    C = 8\pi\epsilon_0 R \ln 2.
    \end{equation*}
    \end{problem}
    \begin{solution}
    \end{solution}
    \begin{note}
    \end{note}

    \item \begin{problem}
    导出圆柱与无限大平板电容:
    \begin{equation*}
    C = \dfrac{2\pi\epsilon_0 L}{\operatorname{arccosh}(d/r)}.
    \end{equation*}
    \end{problem}
    \begin{solution}
    \end{solution}
    \begin{note}
    \end{note}

    \item \begin{problem}
    导出球与无限大平板电容:
    \begin{equation*}
    C = 8\pi\epsilon_0 \sqrt{d^2 - R^2} \sum_{n} \dfrac{R^n}{\left(d+\sqrt{d^2-R^2}\right)^n - \left(d-\sqrt{d^2-R^2}\right)^n}.
    \end{equation*}
    \end{problem}
    \begin{solution}
    \end{solution}
    \begin{note}
    \end{note}

    \item \begin{problem}
    导出两无限长平行圆柱电容:
    \begin{equation*}
    C = \dfrac{2\pi\epsilon_0 L}{\operatorname{arccosh}\left( \dfrac{d^2 - R_1^2 - R_2^2}{2R_1R_2} \right)}, \quad (d > |R_1 - R_2|).
    \end{equation*}
    \end{problem}
    \begin{solution}
    \end{solution}
    \begin{note}
    \end{note}

    \item \begin{problem}
    导出垂直电场中的无限长介质柱极化:
    \begin{equation*}
    \vec{P} = \dfrac{\epsilon_r-1}{\epsilon_r+1} 2\epsilon_0 \vec{E}_0.
    \end{equation*}
    \end{problem}
    \begin{solution}
    \end{solution}
    \begin{note}
    \end{note}

    \item \begin{problem}
    导出均匀外场中的介质球极化:
    \begin{equation*}
    \vec{P} = \dfrac{\epsilon_r-1}{\epsilon_r+2} 3\epsilon_0 \vec{E}_0.
    \end{equation*}
    \end{problem}
    \begin{solution}
    \end{solution}
    \begin{note}
    \end{note}

    \item \begin{problem}
    导出大量均匀介质球的等效极化率:
    \begin{equation*}
    \epsilon_{\text{eff}} = 1 + n \dfrac{\epsilon_r-1}{2\epsilon_r+1} 4\pi R^3.
    \end{equation*}
    \end{problem}
    \begin{solution}
    \end{solution}
    \begin{note}
    \end{note}

    \item \begin{problem}
    导出磁偶极层的磁场:
    \begin{equation*}
    \vec{B} = \dfrac{\mu_0 I}{4\pi} \nabla \Omega.
    \end{equation*}
    \end{problem}
    \begin{solution}
    \end{solution}
    \begin{note}
    \end{note}

    \item \begin{problem}
    证明磁矩在时空缓变电磁场下守恒:
    \begin{equation*}
    \mu = \text{const}.
    \end{equation*}
    \end{problem}
    \begin{solution}
    \end{solution}
    \begin{note}
    \end{note}

    \item \begin{problem}
    证明在 $\vec{B} = (0,0,B)$, $\vec{E} = (0,E,0)$, $\vec{f} = -\eta \vec{v}$, $\vec{F} = -k \vec{r}$, $\vec{r}_0 = (0,0,0)$, $\vec{v}_0 = (v_{x0}, v_{y0}, 0)$ 条件下有:
    \begin{align*}
    \bar{r}(t) &= -\frac{iEq}{k} e^{\left( \frac{iBq+\eta}{2m} \right) t} \frac{ e^{ \sqrt{ \left( \frac{iBq+\eta}{2m} \right)^2 - \frac{k}{m} } t } + e^{ -\sqrt{ \left( \frac{iBq+\eta}{2m} \right)^2 - \frac{k}{m} } t } }{2} \\
               &+ \frac{v_{x0}+iv_{y0} - \frac{iBq+\eta}{2m} \cdot \frac{iEq}{k} }{ \sqrt{ \left( \frac{iBq+\eta}{2m} \right)^2 - \frac{k}{m} } } e^{ \left( \frac{iBq+\eta}{2m} \right) t } \frac{ e^{ \sqrt{ \left( \frac{iBq+\eta}{2m} \right)^2 - \frac{k}{m} } t } - e^{ -\sqrt{ \left( \frac{iBq+\eta}{2m} \right)^2 - \frac{k}{m} } t } }{2} + \frac{iEq}{k}.
    \end{align*}
    \end{problem}
    \begin{solution}
    \end{solution}
    \begin{note}
    \end{note}

    \item \begin{problem}
    证明 Fresnel 公式:
    \begin{equation*}
    \left( \dfrac{E_{0r}}{E_{0i}} \right)_s = \dfrac{n_1 \cos\theta_i - n_2 \cos\theta_t}{n_1 \cos\theta_i + n_2 \cos\theta_t}, \quad \left( \dfrac{E_{0t}}{E_{0i}} \right)_s = \dfrac{2 n_1 \cos\theta_i}{n_1 \cos\theta_i + n_2 \cos\theta_t},
    \end{equation*}
    \begin{equation*}
    \left( \dfrac{E_{0r}}{E_{0i}} \right)_p = \dfrac{n_2 \cos\theta_i - n_1 \cos\theta_t}{n_2 \cos\theta_i + n_1 \cos\theta_t}, \quad \left( \dfrac{E_{0t}}{E_{0i}} \right)_p = \dfrac{2 n_1 \cos\theta_i}{n_2 \cos\theta_i + n_1 \cos\theta_t}.
    \end{equation*}
    \end{problem}
    \begin{solution}
    \end{solution}
    \begin{note}
    \end{note}

    \item \begin{problem}
    导出隐失波沿界面的能流密度:
    \begin{equation*}
    S_x = \dfrac{1}{2} \sqrt{\dfrac{\epsilon_0}{\mu_0}} \dfrac{\sin\theta}{\sin\theta_c} |E_0''|^2 e^{-2\kappa z}.
    \end{equation*}
    \end{problem}
    \begin{solution}
    \end{solution}
    \begin{note}
    \end{note}

    \item \begin{problem}
    导出一维光子晶体 s, p 偏振的色散关系:
    \begin{equation*}
    \cos(K\Lambda) = \cos(k_{az} d_a) \cos(k_{bz} d_b) - \dfrac{1}{2} \left( \dfrac{k_1}{k_2} \dfrac{\mu_b}{\mu_a} + \dfrac{k_2}{k_1} \dfrac{\mu_a}{\mu_b} \right) \sin(k_{az} d_a) \sin(k_{bz} d_b),
    \end{equation*}
    \begin{equation*}
    \cos(K\Lambda) = \cos(k_{az} d_a) \cos(k_{bz} d_b) - \dfrac{1}{2} \left( \dfrac{k_1}{k_2} \dfrac{\epsilon_b}{\epsilon_a} + \dfrac{k_2}{k_1} \dfrac{\epsilon_a}{\epsilon_b} \right) \sin(k_{az} d_a) \sin(k_{bz} d_b).
    \end{equation*}
    \end{problem}
    \begin{solution}
    \end{solution}
    \begin{note}
    \end{note}

    \item \begin{problem}
    导出 Lorentz 模型和 Drude 模型:
    \begin{equation*}
    \epsilon_r(\omega) = 1 + \dfrac{\omega_p^2 (\omega_0^2 - \omega^2)}{(\omega_0^2 - \omega^2)^2 + \gamma^2 \omega^2}, \quad \epsilon_r''(\omega) = \dfrac{\omega_p^2 \gamma \omega}{(\omega_0^2 - \omega^2)^2 + \gamma^2 \omega^2},
    \end{equation*}
    \begin{equation*}
    \epsilon_r(\omega) = 1 - \dfrac{\omega_p^2}{\omega^2 + \gamma^2}, \quad \epsilon_r''(\omega) = \dfrac{\omega_p^2 \gamma}{\omega(\omega^2 + \gamma^2)}.
    \end{equation*}
    \end{problem}
    \begin{solution}
    \end{solution}
    \begin{note}
    \end{note}

    \item \begin{problem}
    导出电磁波在单轴晶体中传播的折射率椭球方程:
    \begin{equation*}
    \dfrac{x^2}{n_0^2} + \dfrac{y^2}{n_0^2} + \dfrac{z^2}{n_e^2} = 1.
    \end{equation*}
    \end{problem}
    \begin{solution}
    \end{solution}
    \begin{note}
    \end{note}

    \item \begin{problem}
    导出在旋光介质中传播的平面波满足的本征方程:
    \begin{equation*}
    \begin{pmatrix} \epsilon & i\epsilon_g & 0 \\ i\epsilon_g & \epsilon & 0 \\ 0 & 0 & \epsilon_z \end{pmatrix} \vec{E} = n^2 \vec{E}.
    \end{equation*}
    \end{problem}
    \begin{solution}
    \end{solution}
    \begin{note}
    \end{note}

    \item \begin{problem}
    导出电磁波在磁化等离子体中传播的介电张量:
    \begin{equation*}
    \bar{\bar{\epsilon}} = \epsilon_0 \begin{pmatrix} \epsilon_1 & i\epsilon_2 & 0 \\ -i\epsilon_2 & \epsilon_1 & 0 \\ 0 & 0 & \epsilon_3 \end{pmatrix},
    \end{equation*}
    其中 $\epsilon_1 = 1 - \dfrac{\omega_p^2}{\omega^2 - \omega_c^2}$, $\epsilon_2 = \dfrac{\omega_c}{\omega} \dfrac{\omega_p^2}{\omega^2 - \omega_c^2}$.
    \end{problem}
    \begin{solution}
    \end{solution}
    \begin{note}
    \end{note}

    \item \begin{problem}
    导出双轴晶体中波矢面方程:
    \begin{equation*}
    \dfrac{k_x^2}{n_1^{-2} - \omega^2/c^2} + \dfrac{k_y^2}{n_2^{-2} - \omega^2/c^2} + \dfrac{k_z^2}{n_3^{-2} - \omega^2/c^2} = 0.
    \end{equation*}
    \end{problem}
    \begin{solution}
    \end{solution}
    \begin{note}
    \end{note}

    \item \begin{problem}
    证明电磁波在旋电介质中传播的色散关系:
    \begin{equation*}
    \det \left( k^2 \delta_{ij} - k_i k_j - \dfrac{\omega^2}{c^2} \mu_{ik} \epsilon_{kj} \right) = 0.
    \end{equation*}
    \end{problem}
    \begin{solution}
    \end{solution}
    \begin{note}
    \end{note}

    \item \begin{problem}
    导出电磁波在双轴晶体中传播的两个可能的折射率:
    \begin{equation*}
    n_{\pm}^2 = \dfrac{2}{ \left( \dfrac{\sin^2\theta}{n_2^2} + \dfrac{\cos^2\theta}{n_3^2} + \dfrac{1}{n_1^2} \right) \pm \sqrt{ \left( \dfrac{\sin^2\theta}{n_2^2} + \dfrac{\cos^2\theta}{n_3^2} + \dfrac{1}{n_1^2} \right)^2 + \dfrac{4\sin^2\theta \cos^2\theta}{n_2^2 n_3^2} } }.
    \end{equation*}
    \end{problem}
    \begin{solution}
    \end{solution}
    \begin{note}
    \end{note}

    \item \begin{problem}
    导出电磁波在磁化铁氧体中传播的导磁张量:
    \begin{equation*}
    \bar{\bar{\mu}} = \mu_0 \begin{pmatrix} \mu & i\kappa & 0 \\ -i\kappa & \mu & 0 \\ 0 & 0 & 1 \end{pmatrix},
    \end{equation*}
    其中 $\mu = 1 + \dfrac{\omega_0 \omega_m}{\omega_0^2 - \omega^2}$, $\kappa = \dfrac{\omega \omega_m}{\omega_0^2 - \omega^2}$.
    \end{problem}
    \begin{solution}
    \end{solution}
    \begin{note}
    \end{note}


    \item \begin{problem}
    导出电磁波在单轴晶体中传播的两个本征偏振态对应的折射率:
    \begin{equation*}
    n_0 \quad \text{和} \quad \tilde{n}_e(\theta) = \left( \dfrac{\cos^2\theta}{n_0^2} + \dfrac{\sin^2\theta}{n_e^2} \right)^{-1/2}.
    \end{equation*}
    \end{problem}
    \begin{solution}
    \end{solution}
    \begin{note}
    \end{note}


    \item \begin{problem}
    导出法拉第旋转角:
    \begin{equation*}
    \theta = \mathcal{V} B d.
    \end{equation*}
    \end{problem}
    \begin{solution}
    \end{solution}
    \begin{note}
    \end{note}

    \item \begin{problem}
    电偶极辐射功率:
    \begin{equation*}
    P = \dfrac{\mu_0 \omega^4 p_0^2}{12\pi c}.
    \end{equation*}
    \end{problem}
    \begin{solution}
    \end{solution}
    \begin{note}
    \end{note}

    \item \begin{problem}
    导出波导截止波数:
    \begin{equation*}
    k_c = \sqrt{ \left( \dfrac{m\pi}{a} \right)^2 + \left( \dfrac{n\pi}{b} \right)^2 }.
    \end{equation*}
    \end{problem}
    \begin{solution}
    \end{solution}
    \begin{note}
    \end{note}

    \item \begin{problem}
    导出导体中衰减常数:
    \begin{equation*}
    \alpha = \omega \sqrt{ \dfrac{\mu\epsilon}{2} \left[ \sqrt{1+\left( \dfrac{\sigma}{\omega\epsilon} \right)^2} - 1 \right] }.
    \end{equation*}
    \end{problem}
    \begin{solution}
    \end{solution}
    \begin{note}
    \end{note}

    \item \begin{problem}
    证明电磁场不变量:
    \begin{equation*}
    \mathcal{S} = \dfrac{1}{2} \left( \dfrac{E^2}{c^2} - B^2 \right).
    \end{equation*}
    \end{problem}
    \begin{solution}
    \end{solution}
    \begin{note}
    \end{note}

    \item \begin{problem}
    导出相对论变换下:
    \begin{equation*}
    E'_\parallel = E_\parallel, \quad B'_\parallel = B_\parallel.
    \end{equation*}
    \end{problem}
    \begin{solution}
    \end{solution}
    \begin{note}
    \end{note}

    \item \begin{problem}
    导出同步辐射功率:
    \begin{equation*}
    P = \dfrac{2}{3} \dfrac{\mu_0 q^2 c \gamma^4 \omega_B^2}{\beta^2}.
    \end{equation*}
    \end{problem}
    \begin{solution}
    \end{solution}
    \begin{note}
    \end{note}

    \item \begin{problem}
    证明辐射阻尼力:
    \begin{equation*}
    \vec{F}_{\text{rad}} = \dfrac{\mu_0 q^2}{6\pi c} \dot{\vec{a}}.
    \end{equation*}
    \end{problem}
    \begin{solution}
    \end{solution}
    \begin{note}
    \end{note}

    \item \begin{problem}
    导出李纳-维谢尔势:
    \begin{equation*}
    A^\mu = \dfrac{\mu_0}{4\pi} \left[ \dfrac{q u^\mu}{u \cdot (x-r)} \right]_{\text{ret}}.
    \end{equation*}
    \end{problem}
    \begin{solution}
    \end{solution}
    \begin{note}
    \end{note}

    \item \begin{problem}
    导出运动电荷电场:
    \begin{equation*}
    \vec{E} = \dfrac{q}{4\pi\epsilon_0} \dfrac{ (1-\beta^2)(\vec{R} - R\vec{\beta}) }{ [R^2 - (\vec{R} \times \vec{\beta})^2]^{3/2} }.
    \end{equation*}
    \end{problem}
    \begin{solution}
    \end{solution}
    \begin{note}
    \end{note}

    \item \begin{problem}
    证明麦克斯韦方程组协变形式:
    \begin{equation*}
    \partial_\mu F^{\mu\nu} = \mu_0 J^\nu.
    \end{equation*}
    \end{problem}
    \begin{solution}
    \end{solution}
    \begin{note}
    \end{note}

    \item \begin{problem}
    推导矩形金属波导 ($a \times b$, $a>b$) 中主模 $TE_{10}$ 的传输功率表达式:
    \begin{equation*}
    P = \dfrac{ab}{4\eta} E_{10}^2 \sqrt{1 - \left( \dfrac{f_c}{f} \right)^2},
    \end{equation*}
    其中 $f_c = \dfrac{c}{2a}$, $\eta = \sqrt{\dfrac{\mu}{\epsilon}}$, $E_{10}$ 为电场幅值.
    \end{problem}
    \begin{solution}
    \end{solution}
    \begin{note}
    \end{note}

    \item \begin{problem}
    导出圆柱形金属波导 (半径 $R$) 中 $TM_{01}$ 模的截止频率:
    \begin{equation*}
    f_c = \dfrac{x_{01}}{2\pi R \sqrt{\mu\epsilon}},
    \end{equation*}
    $x_{01} \approx 2.405$ 为 $J_0(x)=0$ 的第一个根.
    \end{problem}
    \begin{solution}
    \end{solution}
    \begin{note}
    \end{note}

    \item \begin{problem}
    推导对称介质平板波导 (芯层厚 $d$, 折射率 $n_1$, 包层 $n_2$, $n_1>n_2$) 中 $TE_0$ 导模的色散方程:
    \begin{equation*}
    \tan(\kappa d/2) = \dfrac{\gamma}{\kappa},
    \end{equation*}
    其中 $\kappa = \sqrt{k_0^2 n_1^2 - \beta^2}$, $\gamma = \sqrt{\beta^2 - k_0^2 n_2^2}$, $\beta$ 为传播常数.
    \end{problem}
    \begin{solution}
    \end{solution}
    \begin{note}
    \end{note}

    \item \begin{problem}
    导出圆形介质波导 (光纤, 芯径 $a$, 折射率 $n_1$, 包层 $n_2$) 中 $HE_{11}$ 模的特征方程:
    \begin{equation*}
    \left[ \dfrac{J_m'(u)}{u J_m(u)} + \dfrac{K_m'(w)}{w K_m(w)} \right] \left[ \dfrac{J_m'(u)}{u J_m(u)} + \dfrac{n_2^2}{n_1^2} \dfrac{K_m'(w)}{w K_m(w)} \right] = m^2 \left( \dfrac{1}{u^2} + \dfrac{1}{w^2} \right) \left( \dfrac{1}{u^2} + \dfrac{n_2^2}{n_1^2} \dfrac{1}{w^2} \right),
    \end{equation*}
    其中 $u = a \sqrt{k_0^2 n_1^2 - \beta^2}$, $w = a \sqrt{\beta^2 - k_0^2 n_2^2}$.
    \end{problem}
    \begin{solution}
    \end{solution}
    \begin{note}
    \end{note}

    \item \begin{problem}
    推导椭圆金属波导 (长半轴 $a$, 短半轴 $b$) 中主模 ($\mathrm{TE}_{c11}$) 的截止波数近似解:
    \begin{equation*}
    k_c \approx \dfrac{\pi}{2} \sqrt{ \dfrac{1}{a^2} + \dfrac{1}{b^2} }.
    \end{equation*}
    \end{problem}
    \begin{solution}
    \end{solution}
    \begin{note}
    \end{note}

    \item \begin{problem}
    导出脊形金属波导的等效截止波长近似公式:
    \begin{equation*}
    \lambda_c \approx 2 \left( a - t + \dfrac{2b}{\pi} \ln \left[ \dfrac{1}{\sinh \left( \dfrac{\pi s}{2b} \right)} \right] \right),
    \end{equation*}
    其中 $a,b,s,t$ 为脊形几何尺寸.
    \end{problem}
    \begin{solution}
    \end{solution}
    \begin{note}
    \end{note}

    \item \begin{problem}
    推导无耗传输线电报方程 $\dfrac{\partial^2 V}{\partial z^2} = LC \dfrac{\partial^2 V}{\partial t^2}$ 的通解，并给出特性阻抗 $Z_0 = \sqrt{\dfrac{L}{C}}$ 的表达式.
    \end{problem}
    \begin{solution}
    \end{solution}
    \begin{note}
    \end{note}

    \item \begin{problem}
    导出有耗传输线电报方程 $\dfrac{\partial^2 V}{\partial z^2} = (R+j\omega L)(G+j\omega C) V$ ,传播常数 $\gamma = \sqrt{(R+j\omega L)(G+j\omega C)}$ 和特性阻抗 $Z_0 = \sqrt{\dfrac{R+j\omega L}{G+j\omega C}}$.
    \end{problem}
    \begin{solution}
    \end{solution}
    \begin{note}
    \end{note}

    \item \begin{problem}
    证明均匀传输线输入阻抗公式 $Z_{\text{in}}(l) = Z_0 \dfrac{Z_L + jZ_0 \tan(\beta l)}{Z_0 + jZ_L \tan(\beta l)}$, 其中 $\beta = \omega\sqrt{LC}$ 为相位常数.
    \end{problem}
    \begin{solution}
    \end{solution}
    \begin{note}
    \end{note}

    \item \begin{problem}
    计算无限大正方形电阻网格 (每边电阻为 $R$) 中，坐标 $(m,n)$ 与 $(0,0)$ 两点间的等效电阻表达式 $R_{(m,n)}$, 满足二维差分方程 $4\phi_{m,n} - \phi_{m+1,n} - \phi_{m-1,n} - \phi_{m,n+1} - \phi_{m,n-1} = R I \delta_{m,0} \delta_{n,0}$.
    \end{problem}
    \begin{solution}
    \end{solution}
    \begin{note}
    \end{note}

    \item \begin{problem}
    推导无限长一维电阻网络 (相邻结点间电阻为 $r$) 上, 相隔$n$个正方形的结点之间的等效电阻\begin{equation*}
    R=\dfrac{r}{2}\left[n+\dfrac{\sqrt{3}}{3}\left(1+\left(2-\sqrt{3}\right)^n\right)\right] .
    \end{equation*}
    \end{problem}
    \begin{solution}
    \end{solution}
    \begin{note}
    \end{note}

    \item \begin{problem}
    导出在无限大方形电阻网格中, 任意两格点间等效电阻\begin{equation*}
    R=\dfrac{R}{4\pi} \iint_{\omega_{1,2}=0}^{\omega_{1,2}=2\pi} \dfrac{1 - \cos(m\omega_1 + n\omega_2)}{2-\cos(\omega_1)-\cos(\omega_2)} d\omega_1 d\omega_2.    
    \end{equation*}.
    \end{problem}
    \begin{solution}
    \end{solution}
    \begin{note}
    \end{note}

    \item \begin{problem}
    证明n维有限长方体方体电阻网络 (棱边电阻均为 $r$) 中，坐标 $(x_1,x_2,\ldots,x_n)$ 到任意点 $(y_1,y_2,\ldots,y_n)$ 的等效电阻：
    \begin{align*}
    R=& \frac{r}{\prod_{q=1}^{n} m_{q}} \sum_{k_{1}=0}^{m_{1}-1} \sum_{k_{2}=0}^{m_{2}-1} \cdots \sum_{k_{n}=0}^{m_{n}-1} \frac{2^{q-1} \left[ \prod_{q=1}^{n} \cos \frac{2 \pi k_{q}}{m_{q}} \left( x_{q}-\frac{1}{2} \right) - \prod_{q=1}^{n} \cos \frac{2 \pi k_{q}}{m_{q}} \left( y_{q}-\frac{1}{2} \right) \right]}{n - \sum_{q=1}^{n} \cos \frac{2 \pi k_{q}}{m_{q}}} \\
    & \quad \times \left[ \cos \sum_{q=1}^{n} \frac{2 \pi k_{q}}{m_{q}} \left( x_{q}-\frac{1}{2} \right) - \cos \sum_{q=1}^{n} \frac{2 \pi k_{q}}{m_{q}} \left( y_{q}-\frac{1}{2} \right) \right] \\
    = & \frac{r}{\prod_{q=1}^{n} m_{q}} \sum_{k_{1}=0}^{m_{1}-1} \sum_{k_{2}=0}^{m_{2}-1} \cdots \sum_{k_{n}=0}^{m_{n}-1} \frac{2^{q-1} \left[ \prod_{q=1}^{n} \cos \frac{2 \pi k_{q}}{m_{q}} \left( x_{q}-\frac{1}{2} \right) - \prod_{q=1}^{n} \cos \frac{2 \pi k_{q}}{m_{q}} \left( y_{q}-\frac{1}{2} \right) \right]^{2}}{n-\sum_{q=1}^{n} \cos \frac{2 \pi k_{q}}{m_{q}}}.
    \end{align*}
    \end{problem}
    \begin{solution}
    \end{solution}
    \begin{note}
    \end{note}

    \item \begin{problem}
    导出交流 RLC 串联电路的阻抗 $Z = R + j\left( \omega L - \dfrac{1}{\omega C} \right)$ 与谐振频率 $\omega_0 = \dfrac{1}{\sqrt{LC}}$, 并给出电流共振时的幅频特性曲线表达式.
    \end{problem}
    \begin{solution}
    \end{solution}
    \begin{note}
    \end{note}

    \item \begin{problem}
    证明在正弦稳态下, 任意线性无源单口网络的等效阻抗可表示为 $Z(j\omega) = R(\omega) + jX(\omega)$, 其中实部与虚部满足希尔伯特变换关系 (克拉默斯-克勒尼希关系).
    \end{problem}
    \begin{solution}
    \end{solution}
    \begin{note}
    \end{note}

    \item \begin{problem}
    推导三相平衡交流电路在对称负载下的线电压与相电压关系 $V_L = \sqrt{3} V_P e^{j\pi/6}$, 及复功率表达式 $\vec{S} = \sqrt{3} V_L I_L^*$.
    \end{problem}
    \begin{solution}
    \end{solution}
    \begin{note}
    \end{note}

    \item \begin{problem}
    导出包含互感的耦合谐振电路 (电感 $L_1$, $L_2$, 互感 $M$) 的系统方程, 并证明其简正模频率分裂为 $\omega = \omega_0 / \sqrt{1 \pm k}$, 其中耦合系数 $k = M / \sqrt{L_1 L_2}$.
    \end{problem}
    \begin{solution}
    \end{solution}
    \begin{note}
    \end{note}

    \item \begin{problem}
    计算在直流激励下，半无限平面电阻网络 (边界为一条电阻链) 中，从边界点注入电流时,网络内部任意点的电位分布 $\phi(x, y)$.
    \end{problem}
    \begin{solution}
    \end{solution}
    \begin{note}
    \end{note}

    \item \begin{problem}
    推导无耗、无畸变传输线 ($R/L = G/C$) 的波动方程通解，并证明其特性阻抗为纯实数 $Z_0 = \sqrt{\dfrac{L}{C}}$.
    \end{problem}
    \begin{solution}
    \end{solution}
    \begin{note}
    \end{note}

    \item \begin{problem}
    导出周期加载传输线 (每节长度为 $d$, 并联导纳 $Y$) 的色散关系 $\cos(\beta d) = \cos(kd) - \dfrac{j Y Z_0}{2} \sin(kd)$, 其中 $k = \omega\sqrt{LC}$, 并讨论其通带与禁带.
    \end{problem}
    \begin{solution}
    \end{solution}
    \begin{note}
    \end{note}

    \item \begin{problem}
    计算二维无限电阻网络在存在一个缺失电阻 (缺陷) 时,原点与缺陷附近某点之间等效电阻的变化量 $\Delta R$, 用网络格林函数表示。
    \end{problem}
    \begin{solution}
    \end{solution}
    \begin{note}
    \end{note}

    \item \begin{problem}
    证明对于由纯电阻构成的任意网络,任意两结点 $a$, $b$ 间的等效电阻 $R_{ab}$ 可通过关联矩阵与电导矩阵表示为 $R_{ab} = (\mathbf{e}_a - \mathbf{e}_b)^T \mathbf{G}^+ (\mathbf{e}_a - \mathbf{e}_b)$, 其中 $\mathbf{G}^+$ 为电导矩阵的伪逆.
    \end{problem}
    \begin{solution}
    \end{solution}
    \begin{note}
    \end{note}

    \item \begin{problem}
    推导交流电路中,非正弦周期信号 $v(t)$ 的有效值 $V_{\text{rms}} = \sqrt{ \dfrac{1}{T} \int_0^T v^2(t) dt }$ 及其傅里叶级数展开下的各次谐波功率叠加公式.
    \end{problem}
    \begin{solution}
    \end{solution}
    \begin{note}
    \end{note}
\end{enumerate}
% ========== Part 3 结束 ==========

\newpage
% ========== Part 4 光学 ==========
\section*{Part 4 光学}
\begin{enumerate}[label=\textbf{}, leftmargin=*, resume]
    \item \begin{problem}
    证明程函方程:
    \begin{equation*}
    (\nabla S)^2 = n^2(\vec{r}).
    \end{equation*}
    \end{problem}
    \begin{solution}
    \end{solution}
    \begin{note}
    \end{note}

    \item \begin{problem}
    证明光线方程:
    \begin{equation*}
    \dfrac{d}{ds} \left( n \dfrac{d\vec{r}}{ds} \right) = \nabla n.
    \end{equation*}
    \end{problem}
    \begin{solution}
    \end{solution}
    \begin{note}
    \end{note}

    \item \begin{problem}
    证明斯涅尔定律:
    \begin{equation*}
    n_1 \sin\theta_1 = n_2 \sin\theta_2.
    \end{equation*}
    \end{problem}
    \begin{solution}
    \end{solution}
    \begin{note}
    \end{note}

    \item \begin{problem}
    证明理想球面反射镜成像公式:
    \begin{equation*}
    \dfrac{1}{u} + \dfrac{1}{v} = \dfrac{2}{R}.
    \end{equation*}
    \end{problem}
    \begin{solution}
    \end{solution}
    \begin{note}
    \end{note}

    \item \begin{problem}
    证明薄透镜成像公式:
    \begin{equation*}
    \dfrac{1}{f} = \dfrac{1}{u} + \dfrac{1}{v}.
    \end{equation*}
    \end{problem}
    \begin{solution}
    \end{solution}
    \begin{note}
    \end{note}

    \item \begin{problem}
    证明透镜横向放大率:
    \begin{equation*}
    M = -\dfrac{v}{u}.
    \end{equation*}
    \end{problem}
    \begin{solution}
    \end{solution}
    \begin{note}
    \end{note}

    \item \begin{problem}
    证明光在平行平板中产生的横向位移:
    \begin{equation*}
    \Delta = d \sin\theta \left( 1 - \dfrac{\cos\theta}{n \cos\theta'} \right).
    \end{equation*}
    \end{problem}
    \begin{solution}
    \end{solution}
    \begin{note}
    \end{note}

    \item \begin{problem}
    证明在平方律介质 $n^2(x) = n_0^2 (1 - \alpha^2 x^2)$ 中,近轴光线轨迹为
    \begin{equation*}
    x(z) = x_0 \cos(\alpha z) + \dfrac{\theta_0}{\alpha} \sin(\alpha z).
    \end{equation*}
    \end{problem}
    \begin{solution}
    \end{solution}
    \begin{note}
    \end{note}

    \item \begin{problem}
    证明在球对称分布 $n(r) = n_0 / r$ 中,光线轨迹为圆
    \begin{equation*}
    r = R_0 \sec(\theta - \theta_0).
    \end{equation*}
    \end{problem}
    \begin{solution}
    \end{solution}
    \begin{note}
    \end{note}

    \item \begin{problem}
    证明在折射率指数分布 $n(y) = n_0 e^{\beta y}$ 中,光线轨迹为抛物线
    \begin{equation*}
    y(z) = y_0 + \theta_0 z + \dfrac{\beta z^2}{2}.
    \end{equation*}
    \end{problem}
    \begin{solution}
    \end{solution}
    \begin{note}
    \end{note}

    \item \begin{problem}
    证明在轴向渐变光纤 $n^2(r) = n_0^2 [1 - 2\Delta (r/a)^g]$ 中,当 $g=2$ 时子午光线满足
    \begin{equation*}
    \dfrac{d^2 r}{dz^2} + \left( \dfrac{2\Delta}{a^2 n_0^2} \right) r = 0.
    \end{equation*}
    \end{problem}
    \begin{solution}
    \end{solution}
    \begin{note}
    \end{note}

    \item \begin{problem}
    证明平移矩阵：
    \begin{equation*}
    \mathbf{M} = \begin{pmatrix} 1 & 0 \\ \dfrac{d}{n} & 1 \end{pmatrix}.
    \end{equation*}
    \end{problem}
    \begin{solution}
    \end{solution}
    \begin{note}
    \end{note}

    \item \begin{problem}
    证明折射矩阵：
    \begin{equation*}
    \mathbf{M} = \begin{pmatrix} 1 & \dfrac{n_1-n_2}{r} \\ 0 & 1 \end{pmatrix}.
    \end{equation*}
    \end{problem}
    \begin{solution}
    \end{solution}
    \begin{note}
    \end{note}

    \item \begin{problem}
    证明反射矩阵：
    \begin{equation*}
    \mathbf{M} = \begin{pmatrix} 1 & -\dfrac{2n}{r} \\ 0 & 1 \end{pmatrix}.
    \end{equation*}
    \end{problem}
    \begin{solution}
    \end{solution}
    \begin{note}
    \end{note}

    \item \begin{problem}
    证明薄透镜矩阵：
    \begin{equation*}
    \mathbf{M} = \begin{pmatrix} 1 & -\dfrac{1}{f} \\ 0 & 1 \end{pmatrix}.
    \end{equation*}
    \end{problem}
    \begin{solution}
    \end{solution}
    \begin{note}
    \end{note}

    \item \begin{problem}
    证明在周期透镜序列 (焦距 $f$, 间距 $d$) 中,一个周期的传输矩阵为：
    \begin{equation*}
    \mathbf{M} = \begin{pmatrix} 1 - d/f & d \\ -1/f & 1 \end{pmatrix}.
    \end{equation*}
    \end{problem}
    \begin{solution}
    \end{solution}
    \begin{note}
    \end{note}

    \item \begin{problem}
    证明在渐变折射率介质 $\left( n^2(x) = n_0^2 - \alpha x^2 \right)$ 中,长度为 $L$ 的传输矩阵为
    \begin{equation*}
    \mathbf{M} = \begin{pmatrix} \cos(\alpha L) & \sin(\alpha L)/\alpha \\ -\alpha \sin(\alpha L) & \cos(\alpha L) \end{pmatrix}.
    \end{equation*}
    \end{problem}
    \begin{solution}
    \end{solution}
    \begin{note}
    \end{note}

    \item \begin{problem}
    证明在非均匀介质中,光线传输的普遍矩阵形式为：
    \begin{equation*}
    \begin{pmatrix} x_1 \\ \theta_1 \end{pmatrix} = \begin{pmatrix} C(z) & S(z) \\ C'(z) & S'(z) \end{pmatrix} \begin{pmatrix} x_0 \\ \theta_0 \end{pmatrix}.
    \end{equation*}
    \end{problem}
    \begin{solution}
    \end{solution}
    \begin{note}
    \end{note}


    \item \begin{problem}
    证明拉格朗日-亥姆霍兹不变量:
    \begin{equation*}
    L=ny\alpha.
    \end{equation*}
    \end{problem}
    \begin{solution}
    \end{solution}
    \begin{note}
    \end{note}

    \item \begin{problem}
    证明傍轴近似下高斯光束解满足:
    \begin{equation*}
    w(z) = w_0 \sqrt{1 + (z/z_R)^2}.
    \end{equation*}
    \end{problem}
    \begin{solution}
    \end{solution}
    \begin{note}
    \end{note}

    \item \begin{problem}
    证明下述曲面可使宽光束严格成像
    \begin{equation*}
    \left(n_2^2-n_1^2\right)z^2-2\left(n_2-n_1\right)n_2fz+n_2^2y^2=0.
    \end{equation*}
    \end{problem}
    \begin{solution}
    \end{solution}
    \begin{note}
    \end{note}

    \item \begin{problem}
    证明阿贝正弦条件:
    \begin{equation*}
    ny\sin u = \text{const}.
    \end{equation*}
    \end{problem}
    \begin{solution}
    \end{solution}
    \begin{note}
    \end{note}

    \item \begin{problem}
    证明望远镜中:
    \begin{equation*}
    F=\dfrac{f_of_e}{f_o+f_e-d},u_h=-\dfrac{f_od}{f_o+f_e-d},v_h=\dfrac{f_ed}{f_o+f_e-d}.
    \end{equation*}
    并分伽利略望远镜，开普勒望远镜讨论$F$,$u_h$,$v_h$的符号.
    \end{problem}
    \begin{solution}
    \end{solution}
    \begin{note}
    \end{note}

    \item \begin{problem}
    证明显微镜中:
    \begin{equation*}
    M_{eff}=2\text{N.A.}\dfrac{l_0}{D_e}.
    \end{equation*}
    \end{problem}
    \begin{solution}
    \end{solution}
    \begin{note}
    \end{note}

    \item \begin{problem}
    证明薄透镜相位变换因子:
    \begin{equation*}
    t_l = \exp \left( -i \dfrac{k}{2f} (x^2 + y^2) \right).
    \end{equation*}
    \end{problem}
    \begin{solution}
    \end{solution}
    \begin{note}
    \end{note}

    \item \begin{problem}
    证明菲涅尔衍射积分:
    \begin{equation*}
    U(x,y) = \dfrac{e^{ikz}}{i\lambda z} \int U_0(\xi, \eta) e^{\frac{ik}{2z}[(x-\xi)^2+(y-\eta)^2]} d\xi d\eta.
    \end{equation*}
    \end{problem}
    \begin{solution}
    \end{solution}
    \begin{note}
    \end{note}

    \item \begin{problem}
    证明夫琅禾费衍射是孔径的傅里叶变换:
    \begin{equation*}
    U(x,y) \propto \mathcal{F}\{U_0(\xi, \eta)\}.
    \end{equation*}
    \end{problem}
    \begin{solution}
    \end{solution}
    \begin{note}
    \end{note}

    \item \begin{problem}
    证明单缝衍射光强:
    \begin{equation*}
    I(\theta) = I_0 \left[ \dfrac{\sin(\pi a \sin\theta/\lambda)}{\pi a \sin\theta/\lambda} \right]^2.
    \end{equation*}
    \end{problem}
    \begin{solution}
    \end{solution}
    \begin{note}
    \end{note}

    \item \begin{problem}
    证明圆孔衍射艾里斑光强:
    \begin{equation*}
    I(r) = I_0 \left[ \dfrac{2J_1(k a r / z)}{(k a r / z)} \right]^2.
    \end{equation*}
    \end{problem}
    \begin{solution}
    \end{solution}
    \begin{note}
    \end{note}

    \item \begin{problem}
    证明菲涅尔波带片:
    \begin{equation*}
    \dfrac{1}{a}+\dfrac{1}{b}=\left(2n+1\right)\dfrac{\lambda}{\rho_1^2}.
    \end{equation*}
    平面波入射出射光场：
    \begin{equation*}
    U_2\left(r\right) = \dfrac{A}{2}+\dfrac{A}{\pi i}\sum_{n=1}^{\infty} \dfrac{1}{2n+1} \left[ e^{i \left(2n+1\right) \pi r^2/\rho_1^2} - e^{-i\left(2n+1\right) \pi r^2 /\rho_1^2} \right] .    
    \end{equation*}
    \end{problem}
    \begin{solution}
    \end{solution}
    \begin{note}
    \end{note}

    \item \begin{problem}
    证明单缝菲涅尔衍射振幅：
    \begin{equation*}
    U\left(\vec{r}\right)=\dfrac{\left(1-i\right)Ae^{ik\left(z+z_0\right)}}{2\left(z+z_0\right)}\text{exp}\left[\dfrac{ik}{2}\dfrac{\left(x-x_0\right)^2+y^2}{z+z_0}\right]\int_{w_2}^{w_1}e^{i\dfrac{\pi u^2}{2}} du .
    \end{equation*}
    \end{problem}
    \begin{solution}
    \end{solution}
    \begin{note}
    \end{note}
    
    \item \begin{problem}
    证明双光束干涉光强分布:
    \begin{equation*}
    I = 4I_0 \cos^2 \left( \dfrac{\delta}{2} \right),
    \end{equation*}
    其中 $\delta = \dfrac{2\pi}{\lambda} \Delta L$.
    \end{problem}
    \begin{solution}
    \end{solution}
    \begin{note}
    \end{note}

    \item \begin{problem}
    证明法布里-珀罗干涉仪透射峰位满足:
    \begin{equation*}
    2nd \cos\theta = m\lambda.
    \end{equation*}
    \end{problem}
    \begin{solution}
    \end{solution}
    \begin{note}
    \end{note}

    \item \begin{problem}
    证明牛顿环暗环半径:
    \begin{equation*}
    r_m = \sqrt{m \lambda R}.
    \end{equation*}
    \end{problem}
    \begin{solution}
    \end{solution}
    \begin{note}
    \end{note}

    \item \begin{problem}
    证明迈克尔逊干涉仪中,动镜移动 $\Delta d$ 引起条纹移动数
    \begin{equation*}
    \Delta m = \dfrac{2\Delta d}{\lambda}.
    \end{equation*}
    \end{problem}
    \begin{solution}
    \end{solution}
    \begin{note}
    \end{note}

    \item \begin{problem}
    证明多缝光栅光强:
    \begin{equation*}
    I(\theta) = I_0 \left( \dfrac{\sin\alpha}{\alpha} \right)^2 \left( \dfrac{\sin N\beta}{\sin\beta} \right)^2,
    \end{equation*}
    其中 $\alpha = \dfrac{\pi a \sin\theta}{\lambda}$, $\beta = \dfrac{\pi d \sin\theta}{\lambda}$.
    \end{problem}
    \begin{solution}
    \end{solution}
    \begin{note}
    \end{note}

    \item \begin{problem}
    证明光栅方程:
    \begin{equation*}
    d(\sin\theta_m \pm \sin\theta_i) = m\lambda.
    \end{equation*}
    \end{problem}
    \begin{solution}
    \end{solution}
    \begin{note}
    \end{note}

    \item \begin{problem}
    证明光栅角色散:
    \begin{equation*}
    \dfrac{d\theta}{d\lambda} = \dfrac{m}{d \cos\theta_m}.
    \end{equation*}
    \end{problem}
    \begin{solution}
    \end{solution}
    \begin{note}
    \end{note}

    \item \begin{problem}
    证明光栅分辨率:
    \begin{equation*}
    R = \dfrac{\lambda}{\Delta\lambda} = mN.
    \end{equation*}
    \end{problem}
    \begin{solution}
    \end{solution}
    \begin{note}
    \end{note}

    \item \begin{problem}
    证明闪耀光栅条件:
    \begin{equation*}
    2d \sin\theta_B = m\lambda.
    \end{equation*}
    \end{problem}
    \begin{solution}
    \end{solution}
    \begin{note}
    \end{note}

    \item \begin{problem}
    证明正弦振幅光栅的衍射效率:
    \begin{equation*}
    \eta = \left[ \dfrac{J_1(\phi_m/2)}{\phi_m/4} \right]^2.
    \end{equation*}
    \end{problem}
    \begin{solution}
    \end{solution}
    \begin{note}
    \end{note}

    \item \begin{problem}
    证明瑞利判据下,最小可分辨角距离
    \begin{equation*}
    \theta_{\min} = 1.22 \dfrac{\lambda}{D}.
    \end{equation*}
    \end{problem}
    \begin{solution}
    \end{solution}
    \begin{note}
    \end{note}

    \item \begin{problem}
    证明光学传递函数:
    \begin{equation*}
    \mathcal{H}(f_x, f_y) = \mathcal{F}\{h(x,y)\}.
    \end{equation*}
    \end{problem}
    \begin{solution}
    \end{solution}
    \begin{note}
    \end{note}

    \item \begin{problem}
    证明马吕斯定律:
    \begin{equation*}
    I = I_0 \cos^2\theta.
    \end{equation*}
    \end{problem}
    \begin{solution}
    \end{solution}
    \begin{note}
    \end{note}

    \item \begin{problem}
    证明任意偏振态可用斯托克斯参量 $(S_0, S_1, S_2, S_3)$ 描述,且满足
    \begin{equation*}
    S_0^2 \ge S_1^2 + S_2^2 + S_3^2.
    \end{equation*}
    \end{problem}
    \begin{solution}
    \end{solution}
    \begin{note}
    \end{note}

    \item \begin{problem}
    证明线偏振器的穆勒矩阵为:
    \begin{equation*}
    \dfrac{1}{2} \begin{pmatrix} 
    1 & \cos 2\theta & \sin 2\theta & 0 \\
    \cos 2\theta & \cos^2 2\theta & \sin 2\theta \cos 2\theta & 0 \\
    \sin 2\theta & \sin 2\theta \cos 2\theta & \sin^2 2\theta & 0 \\
    0 & 0 & 0 & 0 
    \end{pmatrix}.
    \end{equation*}
    \end{problem}
    \begin{solution}
    \end{solution}
    \begin{note}
    \end{note}

    \item \begin{problem}
    证明波片 (相位延迟 $\delta$,快轴方位角 $0^\circ$) 的穆勒矩阵为
    \begin{equation*}
    \begin{pmatrix} 
    1 & 0 & 0 & 0 \\
    0 & 1 & 0 & 0 \\
    0 & 0 & \cos\delta & \sin\delta \\
    0 & 0 & -\sin\delta & \cos\delta 
    \end{pmatrix}.
    \end{equation*}
    \end{problem}
    \begin{solution}
    \end{solution}
    \begin{note}
    \end{note}

    \item \begin{problem}
    证明部分偏振光偏振度:
    \begin{equation*}
    P = \dfrac{\sqrt{S_1^2 + S_2^2 + S_3^2}}{S_0}.
    \end{equation*}
    \end{problem}
    \begin{solution}
    \end{solution}
    \begin{note}
    \end{note}

    \item \begin{problem}
    证明椭圆偏振光的方位角 $\psi$ 与椭圆率角 $\chi$ 满足:
    \begin{equation*}
    \tan 2\psi = \dfrac{S_2}{S_1}, \quad \sin 2\chi = \dfrac{S_3}{S_0}.
    \end{equation*}
    \end{problem}
    \begin{solution}
    \end{solution}
    \begin{note}
    \end{note}

    \item \begin{problem}
    证明在考虑偏振效应后,薄膜的透射率公式应修正为:
    \begin{equation*}
    T = \dfrac{T_s + T_p}{2}.
    \end{equation*}
    \end{problem}
    \begin{solution}
    \end{solution}
    \begin{note}
    \end{note}

    \item \begin{problem}
    证明朗伯辐射体的辐射亮度 $L_e$ 为常数,其辐射强度角分布满足:
    \begin{equation*}
    I_e(\theta) = I_0 \cos\theta.
    \end{equation*}
    \end{problem}
    \begin{solution}
    \end{solution}
    \begin{note}
    \end{note}

    \item \begin{problem}
    证明点光源在距离 $R$ 处产生的照度 $E_v$ 满足平方反比定律:
    \begin{equation*}
    E_v = \dfrac{I_v}{R^2}.
    \end{equation*}
    \end{problem}
    \begin{solution}
    \end{solution}
    \begin{note}
    \end{note}

    \item \begin{problem}
    证明发光强度的余弦三次方定律:面元 $dA$ 在与其法线成 $\theta$ 角方向上的照度:
    \begin{equation*}
    dE = \dfrac{I_0 \cos^3\theta}{R^2} dA.
    \end{equation*}
    \end{problem}
    \begin{solution}
    \end{solution}
    \begin{note}
    \end{note}

    \item \begin{problem}
    证明准单色光的时间相干长度:
    \begin{equation*}
    L_c = \dfrac{\lambda_0^2}{\Delta\lambda}.
    \end{equation*}
    \end{problem}
    \begin{solution}
    \end{solution}
    \begin{note}
    \end{note}

    \item \begin{problem}
    证明在非相干光照明下,两个空间部分相干点源的干涉条纹可见度:
    \begin{equation*}
    V = |\gamma_{12}(\tau)|.
    \end{equation*}
    \end{problem}
    \begin{solution}
    \end{solution}
    \begin{note}
    \end{note}

    \item \begin{problem}
    证明扩展光源产生的干涉条纹可见度与光源尺寸满足:
    \begin{equation*}
    V = \left| \text{sinc} \left( \dfrac{\pi b \xi}{\lambda z} \right) \right|.
    \end{equation*}
    \end{problem}
    \begin{solution}
    \end{solution}
    \begin{note}
    \end{note}

    \item \begin{problem}
    证明部分相干光的互强度 $J_{12}(\tau)$ 与互相干函数 $\Gamma_{12}(\tau)$ 满足:
    \begin{equation*}
    \Gamma_{12}(\tau) = \langle U_1(t+\tau) U_2^*(t) \rangle.
    \end{equation*}
    \end{problem}
    \begin{solution}
    \end{solution}
    \begin{note}
    \end{note}

    \item \begin{problem}
    证明范西特-泽尼克定理:扩展准单色光源在观察面上产生的复相干度:
    \begin{equation*}
    \mu_{12} = \dfrac{ \iint_S I(\xi,\eta) e^{i k (\Delta R)} d\xi d\eta }{ \iint_S I(\xi,\eta) d\xi d\eta }.
    \end{equation*}
    \end{problem}
    \begin{solution}
    \end{solution}
    \begin{note}
    \end{note}

    \item \begin{problem}
    证明在迈克尔逊测星干涉仪中,条纹可见度与恒星角直径 $\alpha$ 的关系为:
    \begin{equation*}
    V = \left| \dfrac{2J_1 \left( \dfrac{\pi B \alpha}{\lambda} \right) }{ \dfrac{\pi B \alpha}{\lambda} } \right|.
    \end{equation*}
    \end{problem}
    \begin{solution}
    \end{solution}
    \begin{note}
    \end{note}

    \item \begin{problem}
    证明两束部分相干光叠加后的光强:
    \begin{equation*}
    I = I_1 + I_2 + 2\sqrt{I_1 I_2} |\gamma_{12}(\tau)| \cos[\alpha_{12}(\tau)].
    \end{equation*}
    \end{problem}
    \begin{solution}
    \end{solution}
    \begin{note}
    \end{note}

    \item \begin{problem}
    证明高斯-谢尔模型光束的横向相干长度:
    \begin{equation*}
    \rho_0 = \dfrac{\lambda}{\pi \theta_s}.
    \end{equation*}
    \end{problem}
    \begin{solution}
    \end{solution}
    \begin{note}
    \end{note}

    \item \begin{problem}
    证明在二阶非线性过程中,倍频效率:
    \begin{equation*}
    \eta \propto \dfrac{d_{\text{eff}}^2 L^2 I_\omega}{n^3 \lambda^2} \text{sinc}^2 \left( \dfrac{\Delta k L}{2} \right).
    \end{equation*}
    \end{problem}
    \begin{solution}
    \end{solution}
    \begin{note}
    \end{note}

    \item \begin{problem}
    证明光在非线性介质中耦合波方程:
    \begin{equation*}
    \dfrac{dA_1}{dz} = i\kappa A_2 A_3^* e^{i\Delta k z}.
    \end{equation*}
    \end{problem}
    \begin{solution}
    \end{solution}
    \begin{note}
    \end{note}
\end{enumerate}
% ========== Part 4 结束 ==========

\newpage
% ========== Part 5 近代物理 ==========
\section*{Part 5 近代物理}
\begin{enumerate}[label=\textbf{}, leftmargin=*, resume]
    \item \begin{problem}
    证明理想 pn 结在正偏压 $V$ 下的电流密度满足
    \begin{equation*}
    J = J_s \left( e^{eV/(k_B T)} - 1 \right).
    \end{equation*}
    \end{problem}
    \begin{solution}
    \end{solution}
    \begin{note}
    \end{note}

    \item \begin{problem}
    证明两平行理想导体板间的卡西米尔力为
    \begin{equation*}
    F = -\dfrac{\pi^2 \hbar c A}{240 d^4},
    \end{equation*}
    其中 $A$ 为板面积, $d$ 为间距.
    \end{problem}
    \begin{solution}
    \end{solution}
    \begin{note}
    \end{note}

    \item \begin{problem}
    证明在椭圆轨道修正的玻尔模型中,氢原子能级仍为
    \begin{equation*}
    E_n = -\dfrac{m e^4}{2(4\pi\epsilon_0)^2 \hbar^2 n^2}.
    \end{equation*}
    \end{problem}
    \begin{solution}
    \end{solution}
    \begin{note}
    \end{note}

    \item \begin{problem}
    证明考虑相对论修正的玻尔模型能级为
    \begin{equation*}
    E_{n,j} = -\dfrac{m c^2 \alpha^2}{2 n^2} \left[ 1 + \dfrac{\alpha^2}{n} \left( \dfrac{1}{j+\frac{1}{2}} - \dfrac{3}{4n} \right) \right],
    \end{equation*}
    其中 $\alpha$ 为精细结构常数.
    \end{problem}
    \begin{solution}
    \end{solution}
    \begin{note}
    \end{note}

    \item \begin{problem}
    证明康普顿散射中光子波长偏移
    \begin{equation*}
    \Delta\lambda = \dfrac{h}{m_e c} (1 - \cos\theta).
    \end{equation*}
    \end{problem}
    \begin{solution}
    \end{solution}
    \begin{note}
    \end{note}

    \item \begin{problem}
    证明相对论性前灯效应中,辐射角分布满足
    \begin{equation*}
    \dfrac{dP}{d\Omega} \propto \dfrac{1}{\gamma^4 (1-\beta \cos\theta)^4}.
    \end{equation*}
    \end{problem}
    \begin{solution}
    \end{solution}
    \begin{note}
    \end{note}

    \item \begin{problem}
    证明经典霍尔效应中霍尔电压
    \begin{equation*}
    V_H = \dfrac{I B}{n e t},
    \end{equation*}
    其中 $n$ 为载流子浓度, $t$ 为样品厚度.
    \end{problem}
    \begin{solution}
    \end{solution}
    \begin{note}
    \end{note}

    \item \begin{problem}
    证明在均匀磁场 $B$ 中,电子的朗道能级为
    \begin{equation*}
    E_n = \hbar\omega_c \left( n + \dfrac{1}{2} \right),
    \end{equation*}
    其中 $\omega_c = \dfrac{eB}{m}$.
    \end{problem}
    \begin{solution}
    \end{solution}
    \begin{note}
    \end{note}

    \item \begin{problem}
    证明量子霍尔效应中霍尔电阻
    \begin{equation*}
    R_H = \dfrac{h}{\nu e^2},
    \end{equation*}
    $\nu$ 为填充数.
    \end{problem}
    \begin{solution}
    \end{solution}
    \begin{note}
    \end{note}

    \item \begin{problem}
    证明一维无限深方势阱中粒子的能级
    \begin{equation*}
    E_n = \dfrac{n^2 \pi^2 \hbar^2}{2m L^2}.
    \end{equation*}
    \end{problem}
    \begin{solution}
    \end{solution}
    \begin{note}
    \end{note}

    \item \begin{problem}
    证明谐振子能级
    \begin{equation*}
    E_n = \hbar\omega \left( n + \dfrac{1}{2} \right).
    \end{equation*}
    \end{problem}
    \begin{solution}
    \end{solution}
    \begin{note}
    \end{note}

    \item \begin{problem}
    证明氢原子基态波函数
    \begin{equation*}
    \psi_{100} = \dfrac{1}{\sqrt{\pi a_0^3}} e^{-r/a_0},
    \end{equation*}
    其中 $a_0$ 为玻尔半径.
    \end{problem}
    \begin{solution}
    \end{solution}
    \begin{note}
    \end{note}

    \item \begin{problem}
    证明不确定关系
    \begin{equation*}
    \Delta x \Delta p \ge \dfrac{\hbar}{2}.
    \end{equation*}
    \end{problem}
    \begin{solution}
    \end{solution}
    \begin{note}
    \end{note}

    \item \begin{problem}
    证明泡利矩阵满足对易关系
    \begin{equation*}
    [\sigma_i, \sigma_j] = 2i\epsilon_{ijk} \sigma_k.
    \end{equation*}
    \end{problem}
    \begin{solution}
    \end{solution}
    \begin{note}
    \end{note}

    \item \begin{problem}
    证明隧道效应中,方势垒的透射系数
    \begin{equation*}
    T \approx e^{-2a\sqrt{2m(V_0-E)}/\hbar}.
    \end{equation*}
    \end{problem}
    \begin{solution}
    \end{solution}
    \begin{note}
    \end{note}

    \item \begin{problem}
    证明在周期 $\delta$ 势 $V(x) = \alpha \sum_{n=-\infty}^{\infty} \delta(x - n a)$ 中,电子波函数满足的边界条件为
    \begin{equation*}
    \psi'(0^+) - \psi'(0^-) = \dfrac{2m\alpha}{\hbar^2} \psi(0).
    \end{equation*}
    \end{problem}
    \begin{solution}
    \end{solution}
    \begin{note}
    \end{note}

    \item \begin{problem}
    证明克朗尼格-潘纳模型的能带方程
    \begin{equation*}
    \cos(ka) = \cos(qa) + \dfrac{m\alpha}{\hbar^2 q} \sin(qa),
    \end{equation*}
    其中 $q = \sqrt{2mE}/\hbar$.
    \end{problem}
    \begin{solution}
    \end{solution}
    \begin{note}
    \end{note}

    \item \begin{problem}
    证明在紧束缚近似下,一维单原子链的能带为
    \begin{equation*}
    E(k) = E_0 - 2J \cos(ka),
    \end{equation*}
    其中 $J$ 为交叠积分.
    \end{problem}
    \begin{solution}
    \end{solution}
    \begin{note}
    \end{note}

    \item \begin{problem}
    证明在近自由电子近似下,一维晶格在布里渊区边界 $k = \pm \pi/a$ 处出现能隙
    \begin{equation*}
    \Delta = 2|V_G|,
    \end{equation*}
    其中 $V_G$ 为相应倒格矢的傅里叶分量.
    \end{problem}
    \begin{solution}
    \end{solution}
    \begin{note}
    \end{note}

    \item \begin{problem}
    证明布洛赫定理
    \begin{equation*}
    \psi_{n\vec{k}}(\vec{r}+\vec{R}) = e^{i\vec{k}\cdot\vec{R}} \psi_{n\vec{k}}(\vec{r}).
    \end{equation*}
    \end{problem}
    \begin{solution}
    \end{solution}
    \begin{note}
    \end{note}

    \item \begin{problem}
    证明电子有效质量 $m^*$ 满足
    \begin{equation*}
    \dfrac{1}{m^*} = \dfrac{1}{\hbar^2} \dfrac{d^2 E}{dk^2}.
    \end{equation*}
    \end{problem}
    \begin{solution}
    \end{solution}
    \begin{note}
    \end{note}

    \item \begin{problem}
    证明在德拜模型下,固体低温比热
    \begin{equation*}
    C_V \propto T^3.
    \end{equation*}
    \end{problem}
    \begin{solution}
    \end{solution}
    \begin{note}
    \end{note}

    \item \begin{problem}
    证明自由电子的费米能
    \begin{equation*}
    E_F = \dfrac{\hbar^2}{2m} (3\pi^2 n)^{2/3}.
    \end{equation*}
    \end{problem}
    \begin{solution}
    \end{solution}
    \begin{note}
    \end{note}

    \item \begin{problem}
    证明自由粒子的路径积分传播子
    \begin{equation*}
    K(x_b, t_b; x_a, t_a) = \sqrt{ \dfrac{m}{2\pi i\hbar (t_b-t_a)} } \exp\left[ \dfrac{i m (x_b - x_a)^2}{2\hbar (t_b-t_a)} \right].
    \end{equation*}
    \end{problem}
    \begin{solution}
    \end{solution}
    \begin{note}
    \end{note}

    \item \begin{problem}
    证明谐振子的路径积分传播子
    \begin{equation*}
    K(x_b,t_b;x_a,t_a) = \sqrt{ \dfrac{m\omega}{2\pi i\hbar \sin\omega T} } \exp\left\{ \dfrac{im\omega}{2\hbar \sin\omega T} \left[ (x_a^2+x_b^2)\cos\omega T - 2x_a x_b \right] \right\},
    \end{equation*}
    $s.t.\,T=t_b-t_a$.
    \end{problem}
    \begin{solution}
    \end{solution}
    \begin{note}
    \end{note}

    \item \begin{problem}
    证明路径积分表述下的跃迁振幅
    \begin{equation*}
    \langle x_b | e^{-i\hat{H}t/\hbar} | x_a \rangle = \int \mathcal{D}[x(t)] \, e^{\frac{i}{\hbar} S[x(t)]}.
    \end{equation*}
    \end{problem}
    \begin{solution}
    \end{solution}
    \begin{note}
    \end{note}

    \item \begin{problem}
    证明在经典路径附近展开,作用量可写为
    \begin{equation*}
    S[x(t)] = S_{\text{cl}} + \delta^2 S + \cdots.
    \end{equation*}
    \end{problem}
    \begin{solution}
    \end{solution}
    \begin{note}
    \end{note}

    \item \begin{problem}
    证明瞬子解在双势阱隧穿计算中给出的能级分裂
    \begin{equation*}
    \Delta E \propto e^{-S_0/\hbar},
    \end{equation*}
    其中 $S_0$ 为欧氏作用量.
    \end{problem}
    \begin{solution}
    \end{solution}
    \begin{note}
    \end{note}

    \item \begin{problem}
    证明相对论性多普勒效应频率公式
    \begin{equation*}
    \nu' = \nu \sqrt{\dfrac{1-\beta}{1+\beta}},
    \end{equation*}
    其中 $\beta = v/c$, 光源与观察者相背运动.
    \end{problem}
    \begin{solution}
    \end{solution}
    \begin{note}
    \end{note}

    \item \begin{problem}
    证明相对论性粒子的拉格朗日量
    \begin{equation*}
    L = -m_0 c^2 \sqrt{1-\beta^2}.
    \end{equation*}
    \end{problem}
    \begin{solution}
    \end{solution}
    \begin{note}
    \end{note}

    \item \begin{problem}
    证明核反应 $A(a,b)B$ 的反应能
    \begin{equation*}
    Q = (m_A + m_a - m_B - m_b)c^2.
    \end{equation*}
    \end{problem}
    \begin{solution}
    \end{solution}
    \begin{note}
    \end{note}

    \item \begin{problem}
    证明放射性衰变规律
    \begin{equation*}
    N(t) = N_0 e^{-\lambda t},
    \end{equation*}
    其中 $\lambda$ 为衰变常数.
    \end{problem}
    \begin{solution}
    \end{solution}
    \begin{note}
    \end{note}

    \item \begin{problem}
    证明半衰期 $T_{1/2}$ 与衰变常数关系
    \begin{equation*}
    T_{1/2} = \dfrac{\ln 2}{\lambda}.
    \end{equation*}
    \end{problem}
    \begin{solution}
    \end{solution}
    \begin{note}
    \end{note}

    \item \begin{problem}
    证明 $\alpha$ 衰变中,衰变常数 $\lambda$ 与伽莫夫因子 $G$ 满足
    \begin{equation*}
    \lambda \propto e^{-G},
    \end{equation*}
    其中 $G = \dfrac{2\pi Z_1 Z_2 e^2}{\hbar v}$.
    \end{problem}
    \begin{solution}
    \end{solution}
    \begin{note}
    \end{note}

    \item \begin{problem}
    证明核反应阈能
    \begin{equation*}
    E_{\text{th}} = \dfrac{m_A+m_a}{m_A} |Q|,
    \end{equation*}
    当 $Q<0$ 时.
    \end{problem}
    \begin{solution}
    \end{solution}
    \begin{note}
    \end{note}

    \item \begin{problem}
    证明结合能
    \begin{equation*}
    B = [Z m_p + (A-Z)m_n - M_{\text{核}}] c^2.
    \end{equation*}
    \end{problem}
    \begin{solution}
    \end{solution}
    \begin{note}
    \end{note}

    \item \begin{problem}
    证明液滴模型半经验质量公式
    \begin{equation*}
    B(A,Z) = a_V A - a_S A^{2/3} - a_C \dfrac{Z^2}{A^{1/3}} - a_A \dfrac{(A-2Z)^2}{A} + \delta(A,Z).
    \end{equation*}
    \end{problem}
    \begin{solution}
    \end{solution}
    \begin{note}
    \end{note}

    \item \begin{problem}
    证明在液滴模型下,最稳定原子核的 $Z$ 满足
    \begin{equation*}
    \dfrac{Z}{A} = \dfrac{1}{2 + \dfrac{a_C}{2a_A} A^{2/3}}.
    \end{equation*}
    \end{problem}
    \begin{solution}
    \end{solution}
    \begin{note}
    \end{note}

    \item \begin{problem}
    证明壳模型中,核子数 $2,8,20,28,50,82,126$ 为幻数.
    \end{problem}
    \begin{solution}
    \end{solution}
    \begin{note}
    \end{note}

    \item \begin{problem}
    证明集体模型 (几何模型) 中,原子核四极形变参数 $\beta$ 与电四极矩 $Q$ 的关系
    \begin{equation*}
    Q = \dfrac{3}{\sqrt{5\pi}} Z R_0^2 \beta.
    \end{equation*}
    \end{problem}
    \begin{solution}
    \end{solution}
    \begin{note}
    \end{note}

    \item \begin{problem}
    证明费米气体模型下,原子核的能级密度
    \begin{equation*}
    \rho(E) \propto \exp(2\sqrt{aE}),
    \end{equation*}
    其中 $a$ 为能级密度参数.
    \end{problem}
    \begin{solution}
    \end{solution}
    \begin{note}
    \end{note}

    \item \begin{problem}
    证明在独立粒子模型下,原子核磁矩的施密特值
    \begin{equation*}
    \mu = g_j j,
    \end{equation*}
    其中 $g_j$ 为朗德 $g$ 因子.
    \end{problem}
    \begin{solution}
    \end{solution}
    \begin{note}
    \end{note}
\end{enumerate}
% ========== Part 5 结束 ==========

\end{document}    

